%% Elsevier Computer Networks version (elsarticle, two-column 5p layout). Same content as cn_onecolumn.tex.
\documentclass[5p,twocolumn]{elsarticle}

\usepackage{amsmath,amssymb,amsfonts}
\let\proof\relax
\let\endproof\relax
\usepackage{amsthm}
\usepackage{algorithm}
\usepackage{algpseudocode}
\usepackage{graphicx}
\usepackage{booktabs}
\usepackage{xcolor}
\usepackage[T1]{fontenc}
\usepackage{lmodern}
\usepackage{xurl}
\setlength{\emergencystretch}{2em}
\setlength{\textfloatsep}{10pt plus 2pt minus 2pt}
\setlength{\dbltextfloatsep}{10pt plus 2pt minus 2pt}
\setlength{\floatsep}{10pt plus 2pt minus 2pt}
\usepackage{tikz}
\journal{Computer Networks}
\usetikzlibrary{arrows.meta,positioning}

\newtheorem{assumption}{Assumption}
\newtheorem{definition}{Definition}
\newtheorem{lemma}{Lemma}
\newtheorem{proposition}{Proposition}
\newtheorem{corollary}{Corollary}
\newtheorem{theorem}{Theorem}
\theoremstyle{remark}
\newtheorem{remark}{Remark}
\newtheorem{example}{Example}

\newcommand{\E}{\mathbb{E}}
\newcommand{\Prob}{\mathbb{P}}
\newcommand{\ind}[1]{\mathbf{1}\{#1\}}
\newcommand{\A}{\mathcal{A}}
\newcommand{\I}{\mathcal{I}}
\newcommand{\lab}{\ell}
\newcommand{\Dbar}{\bar{\Delta}}
\newcommand{\ee}{\varepsilon}
\newcommand{\xh}{\hat{x}}

\begin{document}

\begin{frontmatter}

\title{Speak When Sure, Ask When Unsure: Confidence-Aware Status Updates from On-Board Classifiers in Aerial IoT}

\author{Sameh Hosny}
\ead{sameh.hosny@aua.am}
\affiliation{organization={Electrical and Computer Engineering Department, American University of Armenia},
            city={Yerevan},
            country={Armenia}}

\begin{abstract}
Internet of Things (IoT) devices classify data on board and send a few-byte label instead of raw data, but a small classifier errs, and wrong labels mislead the receiver. An unmanned aerial vehicle (UAV) monitors a scene for a ground station (GS) and, per slot, stays silent, sends its label, or uploads the frame to a more accurate ground classifier, to minimize the average age of incorrect information (AoII) at the GS under an energy budget and an intermittent link. The UAV never observes the true state, so we formulate a constrained partially observable Markov decision process and show that a belief over the current AoII, updated with each frame's calibrated confidence, is a sufficient statistic. We derive each action's one-step expected AoII in closed form, prove that a label agreeing with the GS is never worth sending, and show that drift-plus-penalty reduces to confidence thresholds ordered silent, raw, label. The policy, CA-DPP, meets the energy budget on every sample path, and comes within 25\% of a lower bound on the optimum when its assumptions hold. On Sentinel-2 flood-mapping traces, it lowers the AoII by 52\% to 75\% relative to pre-identification and age-threshold policies at budgets up to 2\,mJ per slot. Against an AoII-threshold policy, it lowers the AoII by $22\%\pm12\%$ at 0.25\,mJ per slot and up to 49\% when the on-board classifier is weak. When the classifier is accurate and energy is ample, CA-DPP and the AoII-threshold policy perform the same within statistical error.
\end{abstract}

\begin{keyword}
Age of incorrect information \sep on-board inference \sep confidence calibration \sep unmanned aerial vehicle \sep Internet of Things \sep Lyapunov optimization \sep POMDP
\end{keyword}

\end{frontmatter}

%% =====================================================================
\section{Introduction}\label{sec:intro}
%% =====================================================================
Many Internet of Things (IoT) devices no longer send what they sense. A camera trap reports that it saw an animal, a nanosatellite reports that a pixel block contains a ship, and an unmanned aerial vehicle (UAV) that surveys a flood reports that the area below it is under water. Each of these devices runs a small classifier on board and transmits its decision, a label of a few bytes, instead of an image of tens of kilobytes~\cite{shamsoshoara2021flame,kyrkou2020aider}. For a battery-powered device on a weak wireless link, this saves most of the energy and airtime that a raw upload would cost.

The saving comes at a price. A classifier small enough to run on a UAV or a microcontroller makes mistakes, and it makes them most often on ambiguous inputs: a flooded field under cloud shadow, a dry field after rain. Once the transmitted content is an inference, a delivered update can be wrong. The receiver then holds a fresh but wrong estimate, and it keeps it until a later update corrects it. The age of information (AoI) does not register this error, because it resets on every delivery whether the content is right or wrong. The age of incorrect information (AoII)~\cite{maatouk2020aoii} grows while the receiver's estimate differs from the true state, so it penalizes wrong estimates as well as stale ones.

Existing AoII work assumes, however, that the transmitter observes the source through a sensor~\cite{maatouk2020aoii,zakeri2023aoiieh,ngo2025eh}. Such a transmitter knows whether the receiver is wrong and can decide to update exactly when it is. A device whose sensor is a classifier has no such knowledge. It sees only its own label and how confident the classifier is about it. A confident label that disagrees with the receiver probably means that the receiver is wrong. An unsure label could mean either, and sending it may turn a correct estimate into a wrong one. The device has a third option: it can upload the raw data and let a larger classifier on the ground decide, at a much higher energy cost that depends on the link.

Consider a UAV that maps a flood for a ground station (GS) that coordinates rescue teams. If a frame looks clearly flooded and the GS status says dry, the UAV should send its label. If the frame is ambiguous, it should upload the frame when the link is good enough, and wait otherwise. If its label agrees with the GS status, it should stay silent. This paper turns these rules of thumb into a policy with guarantees. Four difficulties stand in the way. First, the UAV never observes whether the GS is wrong. Second, the confidence of a small classifier is itself noisy, and its errors are correlated across consecutive frames because consecutive footprints overlap. Third, the delivery probability and energy of each action depend on the distance to the GS and on blockage. Fourth, the UAV has a fixed energy budget for the whole mission, so a transmission now reduces what is available later.

We address three questions. \emph{(Q1)} In each slot, should the UAV stay silent, send its label or upload the frame, given its label, confidence, link and energy? \emph{(Q2)} Does the per-frame confidence lower the AoII, and when? \emph{(Q3)} When is a raw upload worth its cost? We answer Q1 analytically. The UAV keeps a belief over the current AoII of the GS status, updated with the calibrated confidence of each frame and with the acknowledgments of the GS. From the belief we compute in closed form how much each action reduces the expected AoII of the next slot, and a Lyapunov drift-plus-penalty rule weighs this reduction against the energy cost, with a price that rises when the UAV overspends. We answer Q2 and Q3 on real inference traces. The main contributions are as follows.
\begin{enumerate}
\item \emph{Model.} We formulate confidence-aware status updating as a constrained partially observable Markov decision process (POMDP). It couples a calibrated per-sample error model of the on-board classifier, a ground classifier, an intermittent air-to-ground link and an energy budget. We show that a belief over the current AoII is a sufficient statistic, and we derive its update, including a prior correction and a tempered likelihood for errors that are correlated in time (Sections~\ref{sec:model} and~\ref{sec:belief}).
\item \emph{Structure.} We derive the one-step expected AoII of each action in closed form. We prove that a label that agrees with the GS has no value and that the per-slot decision has a threshold structure in the confidence: the UAV stays silent at low confidence, uploads the raw frame at intermediate confidence, and sends the label at high confidence (Section~\ref{sec:structure}).
\item \emph{Policy.} We build the online policy CA-DPP on these results and prove that it meets the energy budget on every sample path. CA-DPP needs no knowledge of the confidence distribution, the blockage statistics or the trajectory, and its per-slot cost is linear in the size of the belief (Section~\ref{sec:lyap}).
\item \emph{Evaluation.} We evaluate CA-DPP on inference traces from the Sen1Floods11 dataset~\cite{bonafilia2020sen1floods11} under chip-level cross-validation, against policies adapted from four recent designs~\cite{zakeri2023aoiieh,gan2025taoi,pehlivanoglu2026orbit,basnayaka2025aomi}. The per-frame confidence pays off under tight energy budgets and with weak on-board classifiers, and raw uploads pay off near the GS at larger budgets; with an accurate classifier and ample energy, simpler policies perform as well. In a synthetic model that satisfies our assumptions, a full-information lower bound places CA-DPP within $14\%$ to $25\%$ of the optimum, and correlated errors show where its advantage fades (Sections~\ref{sec:results} and~\ref{sec:discussion}).
\end{enumerate}

The rest of the paper is organized as follows. Section~\ref{sec:related} reviews related work, Sections~\ref{sec:model} to~\ref{sec:lyap} answer Q1, Section~\ref{sec:results} answers Q2 and Q3, Section~\ref{sec:discussion} discusses the results, Section~\ref{sec:concl} concludes; proofs are in the appendix.
%% =====================================================================
\section{Related Work}\label{sec:related}
%% =====================================================================
We review AoII, the freshness of inferred information, and inference offloading from IoT devices and UAVs.

\subsection{Age of Incorrect Information}
The AoI treats every update as correct, so it cannot tell a fresh wrong estimate from a fresh correct one. Maatouk et al.~\cite{maatouk2020aoii} introduced the AoII to fix this. The AoII is zero while the receiver's estimate equals the source state and grows with time while the two differ. For a Markov source observed perfectly by the transmitter, they showed that the optimal policy transmits only when the estimate is wrong.

Two recent works add energy constraints. Zakeri et al.~\cite{zakeri2023aoiieh} study an energy-harvesting sensor that decides when to sample and when to transmit. Since the sensor does not see the source between samples, they formulate a POMDP whose belief is about the source state, and they compute the AoII-optimal policy numerically, finding a multi-threshold structure. Ngo et al.~\cite{ngo2025eh} study the age and the correctness of updates from many energy-harvesting devices that share a random-access channel. In both works, a sample reveals the source state without error. Our UAV, in contrast, never observes the state. Its only observation is a classifier output that may be wrong, and the probability that it is wrong changes from frame to frame.

The belief we use is closest to that of Cosandal et al.~\cite{cosandal2024joint}, who keep a joint belief over the age and the source state at a monitor that pulls samples with a delay, and use it for estimation and for sampling decisions. There, too, a delivered sample is exact, and the uncertainty comes from the delay. In our problem a fresh observation arrives in every slot, but it may be wrong, so the belief must weigh each frame by how much its confidence can be trusted.

\subsection{Freshness of Inferred Information}
Several works consider status updates whose content is the output of a classifier. Buratto et al.~\cite{buratto2023misclass} add a penalty for misclassification to an AoI scheduling problem. The classifier errs with a fixed probability, and the schedule is computed offline. Gan et al.~\cite{gan2025taoi} define a task-oriented AoI for cameras that pre-identify each frame on the device and send the frames that appear to contain the target to a server, which re-identifies them with a larger model. The error rates of both classifiers are fixed, so the device does not use the confidence of individual frames. Pehlivano\u{g}lu, Li and Uysal~\cite{pehlivanoglu2026orbit} let a satellite choose among idling, running inference on board, sending a cached result and offloading raw data to the ground. Their action set is the closest to ours. They prove that a threshold policy minimizes the age of inference, but each action succeeds with a fixed probability, errors do not depend on the sample, and energy is not constrained. Basnayaka et al.~\cite{basnayaka2025aomi} define the age of misclassified information for event detection from low Earth orbit, again with a constant classification accuracy.

These works establish that the correctness of an inferred update matters for freshness. All of them, however, describe the classifier by a single error rate. A calibrated classifier gives more information than that: it reports, for each frame, how likely its label is to be wrong. We show how to use this information in the decision to stay silent, send a label or ask the ground.

\subsection{Inference Offloading from IoT Devices and UAVs}
A large body of work decides whether an IoT device or a UAV should run a task locally or offload it to an edge server, usually to minimize delay or energy. Yang et al.~\cite{yang2020uav}, for example, split a deep network for UAV target tracking between the UAV and an edge server so as to minimize a weighted sum of delay and energy. Such work accounts for delay and energy, but it does not measure how fresh or how correct the receiver's knowledge is. Datasets such as FLAME~\cite{shamsoshoara2021flame} for wildfires and AIDER~\cite{kyrkou2020aider} for disaster scenes show that classifiers small enough for a UAV reach useful accuracy, but they also show that such classifiers err on a noticeable fraction of the frames. Post-hoc calibration~\cite{guo2017calibration} makes the output score of such a classifier a reliable probability, which is what our belief update requires.

The energy constraint in our problem is handled with Lyapunov drift-plus-penalty~\cite{neely2010book}. Huang and Neely~\cite{huang2013eh} extend this method to devices that harvest energy into a finite battery. Since the AoII depends on past decisions, the standard performance bound of drift-plus-penalty does not apply directly. The renewal framework of~\cite{neely2013renewal} is the natural tool for such problems (Section~\ref{sec:lyap}).

\subsection{Positioning}
Table~\ref{tab:positioning} compares the closest works on four elements: (E1) a per-sample error model driven by the classifier confidence; (E2) AoII on the inferred label; (E3) an energy constraint handled online; and (E4) a choice among silence, a label and raw data. Each prior work covers at most two of these elements fully. Ours combines all four.

\begin{table}[t]
\caption{Positioning against the closest works. \checkmark: covered; $\circ$: partly covered; --: not covered.}
\label{tab:positioning}
\centering\footnotesize
\setlength{\tabcolsep}{3pt}
\resizebox{\columnwidth}{!}{%
\begin{tabular}{@{}lcccc@{}}
\toprule
Work & (E1) per-sample & (E2) AoII on & (E3) online & (E4) silent/\\
 & errors & inferred label & energy & label/raw\\
\midrule
Zakeri et al.~\cite{zakeri2023aoiieh} & -- & \checkmark & $\circ$ & $\circ$\\
Ngo et al.~\cite{ngo2025eh} & -- & $\circ$ & $\circ$ & --\\
Buratto et al.~\cite{buratto2023misclass} & $\circ$ & \checkmark & -- & --\\
Gan et al.~\cite{gan2025taoi} & $\circ$ & $\circ$ & -- & $\circ$\\
Pehlivano\u{g}lu et al.~\cite{pehlivanoglu2026orbit} & -- & $\circ$ & -- & \checkmark\\
Basnayaka et al.~\cite{basnayaka2025aomi} & $\circ$ & $\circ$ & -- & --\\
Yang et al.~\cite{yang2020uav} & -- & -- & $\circ$ & $\circ$\\
\midrule
This work & \checkmark & \checkmark & \checkmark & \checkmark\\
\bottomrule
\end{tabular}}
\end{table}
%% =====================================================================
\section{System Model}\label{sec:model}
%% =====================================================================
\subsection{Network and Scenario}\label{sec:story}
We consider a UAV that flies over a disaster area and reports to a ground station (GS) that coordinates the response, as shown in Fig.~\ref{fig:system}. Time is divided into slots of length $\tau$, indexed by $t\in\{0,1,\dots\}$. The UAV flies at constant speed along a route planned before the mission, so its position $\mathbf{u}(t)$ in every slot is known in advance. In each slot the UAV photographs the patch of ground below it, which we call its footprint, and a small classifier on board decides whether an event of interest is present in the patch, for example flooding, fire, a collapsed building or a survivor. Throughout the paper we use flood mapping as the running example: the classifier labels each photo ``flooded'' or ``dry'', and the response teams need to know which patches are under water. The model applies to any such event with a suitable classifier.

Since the UAV moves, the patch below it changes from slot to slot, and consecutive footprints partly overlap. We do not model the whole area. We model only what the UAV sees along its route: in slot $t$, $X(t)$ is the state of the patch below the UAV and $Y(t)$ is the photo of that patch. The position $\mathbf{u}(t)$ affects the model in one place only, the link to the GS. The distance and the elevation angle to the GS determine how likely a packet is to arrive and how much energy it costs. The route can have any shape. In our evaluation it is a circle, so the UAV passes alternately through parts of the route close to the GS and far from it.

The GS has no sensor of its own. Everything it knows about the patch below the UAV comes from the UAV. We call the GS status the label that the GS currently holds for this patch: it is the last label the GS received, either a label sent by the UAV or the label that a larger classifier at the GS produced from a photo sent by the UAV. The GS holds this status until the next update arrives, and the response teams act on it. In each slot the UAV can therefore stay silent, send its on-board label, or send the photo. A label is a few bytes and cheap to send, but it may be wrong. A photo is tens of kilobytes and much more expensive, and it may fail to arrive far from the GS, but the ground classifier labels it almost always correctly. The UAV chooses under an energy budget and without observing the true state; this last fact shapes the whole analysis. Table~\ref{tab:notation} summarizes the notation.

\begin{figure}[t]
\centering
\includegraphics[width=0.8\columnwidth]{figures/fig_scenario.pdf}
\caption{System model. A UAV flies over a disaster area, photographs the patch below it and labels each photo on board with a confidence $c(t)$. In each slot it stays silent, sends its label (blue) or uploads the photo to the more accurate ground classifier (orange). The GS keeps the last label received and acks each decoded packet. Inset: the AoII grows while the GS status is wrong and resets when it becomes right.}
\label{fig:system}
\end{figure}

\begin{table}[t]
\caption{Main notation}\label{tab:notation}
\centering\footnotesize
\setlength{\tabcolsep}{3pt}
\begin{tabular}{@{}p{0.58\columnwidth}ll@{}}
\toprule
Meaning & Symbol & Defined in\\
\midrule
Patch below the UAV: event or not & $X(t)$ & \eqref{eq:source}\\
How often the scene changes & $p_0,p_1$; $\varpi$ & \eqref{eq:source}, \eqref{eq:stationary}\\
Photo; on-board label, confidence & $Y(t)$; $\lab(t)$, $c(t)$ & \eqref{eq:calib}\\
Label disagrees with the GS status & $d(t)$ & Sec.~\ref{sec:update}\\
Label from the ground classifier & $\lab^{\rm g}(t)$, $\ee_{\rm g}(c)$ & Assum.~\ref{as:inf}\\
Silent, send label, send photo & $a(t)\in\{0,1,2\}$ & Sec.~\ref{sec:act}\\
UAV position, blockage & $g(t)$ & \eqref{eq:block}, \eqref{eq:pl}\\
Packet decoded & $o(t)$ & \eqref{eq:est}\\
Delivery probability; energy & $q_a(g)$; $E_a(g)$ & \eqref{eq:succ}, \eqref{eq:energy}\\
Energy budget of the mission & $\bar E$ & \eqref{eq:budget}\\
What the GS status says & $\hat X(t)$ & \eqref{eq:est}\\
How long the GS status has been wrong & $\Delta(t)$ (AoII) & \eqref{eq:aoii}\\
What the UAV believes about $\Delta(t)$ & $\pi_t$; $\mu_t$, $S_t$ & \eqref{eq:belief}, \eqref{eq:musum}\\
Evidence of a frame that the GS is wrong & $\Lambda_t$ & \eqref{eq:Lambda}\\
Value of an action; its energy price & $G_a$; $Z(t)$, $V$ & Secs.~\ref{sec:structure}, \ref{sec:lyap}\\
\bottomrule
\end{tabular}
\end{table}

\subsection{Event Dynamics}\label{sec:src}
\begin{figure}[t]
\centering
\includegraphics[width=0.55\columnwidth]{figures/fig_markov.pdf}
\caption{Two-state Markov chain of the patch state $X(t)$. A dry patch ($0$) is followed by a flooded one with probability $p_0$, and a flooded patch ($1$) by a dry one with probability $p_1$. With $p_0,p_1<\frac12$ the state tends to persist.}
\label{fig:markov}
\end{figure}
We describe the patch below the UAV by a binary state $X(t)\in\{0,1\}$, with $X(t)=1$ if the event is present and $X(t)=0$ otherwise. In the flood example, $X(t)=1$ means that the patch is flooded. The UAV moves a short distance between slots, so consecutive frames cover overlapping areas and the state tends to persist: a flooded footprint is usually followed by another flooded footprint. Note that $X(t)$ refers to the patch below the UAV in slot $t$, not to a fixed patch. The chain therefore describes how the state changes from one footprint to the next along the route, and its transition probabilities depend on the speed of the UAV and on the overlap of its footprints. They are constant because the speed is. The simplest model that captures this persistence is the two-state Markov chain of Fig.~\ref{fig:markov}. Let $p_x$ be the probability that the state leaves $x$ in one slot,
\begin{equation}\label{eq:source}
\Prob\big\{X(t+1)\neq x\mid X(t)=x\big\}=p_x,\qquad x\in\{0,1\}.
\end{equation}
We assume $0<p_0,p_1<\tfrac12$. The lower bound makes the chain irreducible, so the scene does change and every state is visited. The upper bound says that the state is more likely to stay than to change in one slot. This persistence is what makes an old status useful at all: if the state changed more often than not, the GS would do better by inverting its last label than by keeping it. We use the bound $p_x<\frac12$ in Section~\ref{sec:structure}.

The fraction of time the chain spends in each state follows from the balance between the two transitions. In steady state the probability flow from 0 to 1 equals the flow from 1 to 0, so $\varpi_0p_0=\varpi_1p_1$, and with $\varpi_0+\varpi_1=1$,
\begin{equation}\label{eq:stationary}
\varpi_1=\frac{p_0}{p_0+p_1},\qquad \varpi_0=\frac{p_1}{p_0+p_1}.
\end{equation}
We call $\varpi=(\varpi_0,\varpi_1)$ the class prior. The chain may be asymmetric. In flood mapping, for instance, a flooded footprint typically turns dry more often than a dry one floods ($p_1>p_0$), so $\varpi_1<\varpi_0$. We use the symmetric case $p_0=p_1$ only to state some results in a simpler form.

\subsection{On-Board and Ground Classification}
In slot $t$ the UAV takes a photo $Y(t)$ of the patch below it, which we also call a frame. The photo is used only to classify the patch; the UAV knows its position from the route and does not infer it from the photo. The on-board classifier returns a label and a confidence,
\begin{equation}\label{eq:label}
\lab(t)\in\{0,1\},\qquad c(t)\in[\tfrac12,1],
\end{equation}
where $\lab(t)$ is the state that the classifier finds more likely and $c(t)$ is the probability it assigns to that state. A classifier that outputs a probability $p$ of flooding, as most do, gives $\lab(t)=\ind{p\ge\frac12}$ and $c(t)=\max\{p,1-p\}$. The confidence is useful only if it means what it says. We therefore require that the classifier be calibrated: among all frames that receive label $\lab$ and confidence $c$, a fraction $c$ is truly in state $\lab$, for both label values,
\begin{equation}\label{eq:calib}
\Prob\big\{X(t)=\lab(t)\mid\lab(t),\,c(t)\big\}=c(t).
\end{equation}
Averaging~\eqref{eq:calib} over the two label values shows that a label with confidence $c$ is wrong with probability $1-c$. The converse does not hold. A classifier can have this property while being overconfident when it says ``flooded'' and underconfident when it says ``dry''. We need~\eqref{eq:calib} for each label value because the belief update below uses the probability of each state given the label, and flooded and dry patches are not equally common.

Modern classifiers are often overconfident, so~\eqref{eq:calib} does not hold out of the box. Post-hoc calibration~\cite{guo2017calibration}, such as isotonic regression on held-out data, restores it to a good approximation. Fig.~\ref{fig:inference}(a) shows that our on-board model, calibrated in this way, is close to calibrated (Section~\ref{sec:results} gives the numbers). Fig.~\ref{fig:inference}(b) confirms the consequence above: the measured error rate of the on-board label follows $1-c$ closely across confidence bins.

The belief update in Section~\ref{sec:update} needs the likelihood of a frame under each state, that is, how much more probable the observed pair $(\lab(t),c(t))$ is if $X(t)=1$ than if $X(t)=0$. Calibration gives the posterior probability of each state, not the likelihood. By~\eqref{eq:calib}, the posterior probability of state $x\in\{0,1\}$ is
\begin{equation}\label{eq:cx}
c_x(t)=\Prob\big\{X(t)=x\mid\lab(t),c(t)\big\}=\begin{cases}c(t), & x=\lab(t),\\ 1-c(t), & x\neq\lab(t).\end{cases}
\end{equation}
By Bayes' rule and~\eqref{eq:cx}, posterior and likelihood are related through the class prior,
\begin{equation}\label{eq:lik}
\Prob\big\{\lab(t),c(t)\mid X(t)=x\big\}=\frac{c_x(t)\,\Prob\big\{\lab(t),c(t)\big\}}{\varpi_x}\propto\frac{c_x(t)}{\varpi_x},
\end{equation}
since the factor $\Prob\{\lab(t),c(t)\}$ does not depend on $x$. The division by $\varpi_x$ matters when the classes are unbalanced. A frame with confidence $c(t)=\frac12$, for example, looks neutral, but by~\eqref{eq:lik} it is $\varpi_0/\varpi_1$ times more likely under $X=1$ than under $X=0$, which exceeds one when flooded footprints are the rarer class. The classifier has already used the prior to produce its output, and the belief update below uses the source dynamics, which contain the same prior. Dividing by $\varpi_x$ prevents counting the prior twice.

After a raw upload, the ground classifier labels the frame. We denote its label by $\lab^{\rm g}(t)$ and its error probability by $\ee_{\rm g}$. The ground model is larger and sees all spectral bands, so it is more accurate than the on-board model. Its errors are not independent of the on-board confidence, however: frames that are hard on board are often hard on the ground as well. In our data the ground model errs far more often on the frames that the on-board model gets wrong than on the others, and Fig.~\ref{fig:inference}(b) shows that its error rate falls as the on-board confidence grows. We therefore let the ground error depend on the on-board confidence and write $\ee_{\rm g}(c)$.

\begin{figure*}[t]
\centering
\includegraphics[width=\textwidth]{figures/fig_inference.pdf}
\caption{Classifier outputs on the Sen1Floods11 test chips. (a) Reliability diagram: fraction of truly flooded footprints versus the predicted flood probability, for the on-board model, its noisiest variant and the ground model, with their expected calibration errors (ECE). All three are well calibrated. (b) Error rate of the on-board and ground labels versus the on-board confidence $c$. The on-board error follows $1-c$ (dotted); the ground model also errs more on hard frames. (c) $120$ consecutive footprints of one test chip: flooded footprints shaded, on-board flood probability (line) and wrong labels (crosses). Errors come in runs, mostly missed floods.}
\label{fig:inference}
\end{figure*}

We collect these modeling choices in an assumption. It is used only to design the controller; the simulations in Section~\ref{sec:results} replay real classifier outputs, which only approximate it.

\begin{assumption}[Inference model]\label{as:inf}
(i) The on-board classifier is calibrated, i.e.,~\eqref{eq:calib} holds. (ii) Given the source path, the on-board outputs $(\lab(t),c(t))$ are independent across slots. (iii) Given $X(t)$ and $c(t)=c$, the ground label is wrong with probability $\ee_{\rm g}(c)<\tfrac12$, independently of $\lab(t)$ and of the past.
\end{assumption}

Part (ii) is the strongest of the three. Consecutive footprints overlap, so a feature that fools the classifier in one frame tends to fool it in the next. Fig.~\ref{fig:inference}(c) shows such runs of errors, and Section~\ref{sec:results} measures them: a wrong label is followed by another wrong label far more often than the average error rate would suggest. We keep part (ii) because it yields a tractable belief update, and we correct for its optimism in Section~\ref{sec:update}.

\subsection{Actions and Channel}\label{sec:act}
At the start of slot $t$, after classifying the frame, the UAV chooses an action $a(t)\in\A=\{0,1,2\}$. It stays silent ($a=0$), which costs nothing and leaves the GS status unchanged. It sends its label in a packet of $L_1$ bits ($a=1$), which corrects a wrong GS status if the label is right but makes a correct status wrong if the label is wrong. Or it uploads the frame in a packet of $L_2\gg L_1$ bits ($a=2$), so that the ground classifier labels it. Whether a packet arrives, and how much energy it costs, depends on the link, which we model next.

Whether a packet arrives depends on the link state $g(t)=(\mathbf{u}(t),\beta(t))$. The first part is the position $\mathbf{u}(t)$ of the UAV on its route, which is known in advance because the route and the speed are fixed. The second part is a blockage indicator $\beta(t)\in\{0,1\}$. Buildings and terrain block the link for several slots at a time, so we model $\beta(t)$ as a two-state Markov chain with
\begin{equation}\label{eq:block}
\begin{split}
\Prob\big\{\beta(t+1)=1\mid\beta(t)=0\big\}&=\rho_{\rm b},\\
\Prob\big\{\beta(t+1)=0\mid\beta(t)=1\big\}&=\rho_{\rm u}.
\end{split}
\end{equation}
A blockage then lasts $1/\rho_{\rm u}$ slots on average, and a fraction $\rho_{\rm b}/(\rho_{\rm b}+\rho_{\rm u})$ of the slots is blocked. No packet gets through a blocked link.

When the link is not blocked, the received power depends on the distance $r$ between the UAV and the GS and on the elevation angle $\phi$. A low UAV is often seen through buildings, so the link is in line of sight (LoS) only with some probability. Following~\cite{alhourani2014lap}, we use the mean path loss
\begin{equation}\label{eq:pl}
\overline{\rm PL}(r,\phi)={\rm FSPL}(r)+P_{\rm LoS}(\phi)\,\eta_{\rm LoS}+\big(1-P_{\rm LoS}(\phi)\big)\eta_{\rm NLoS},
\end{equation}
where ${\rm FSPL}(r)$ is the free-space path loss, $\eta_{\rm LoS}$ and $\eta_{\rm NLoS}$ are the excess losses in LoS and non-LoS, and $P_{\rm LoS}(\phi)=[1+\alpha e^{-b(\phi-\alpha)}]^{-1}$ grows with the elevation angle, with environment constants $\alpha$ and $b$. The transmit power, the path loss~\eqref{eq:pl}, the bandwidth $W$ and the receiver noise figure give the mean signal-to-noise ratio (SNR) $\bar\gamma(g)$.

Small-scale fading makes the instantaneous SNR random around this mean. With Rayleigh fading, the instantaneous SNR is $\bar\gamma(g)\,|h|^2$ with $|h|^2$ exponentially distributed with unit mean. The fading coefficient $h$ stays constant during a slot. The UAV sends a packet of $L_a$ bits with a rateless (incremental-redundancy) code: it keeps sending coded bits until the GS decodes the packet and acknowledges it, or until a transmission window $\tau_{\max}\le\tau$ ends. Decoding takes the time $L_a/C(h)$, where $C(h)=W\log_2(1+\bar\gamma|h|^2)$ is the capacity of the faded channel, so the packet is decoded if $L_a/C(h)\le\tau_{\max}$, that is, if $W\tau_{\max}\log_2(1+\bar\gamma|h|^2)\ge L_a$. Solving for $|h|^2$ gives the delivery probability
\begin{equation}\label{eq:succ}
q_a(g)=(1-\beta)\,\exp\!\Big(-\frac{\gamma_a}{\bar\gamma(g)}\Big),\qquad \gamma_a=2^{L_a/(W\tau_{\max})}-1,
\end{equation}
for $a\in\{1,2\}$,
and $q_0\equiv0$. Here $\gamma_a$ is the SNR needed to deliver $L_a$ bits within the window. Since $L_2>L_1$, $\gamma_2>\gamma_1$, so $q_2(g)\le q_1(g)$ in every link state.

The link is therefore an erasure channel: a packet either arrives intact or is lost, and a lost packet leaves the GS status unchanged. The link never changes the content of a label. Every packet carries an $n$-bit cyclic redundancy check (CRC), and the GS discards a packet whose check fails; a corrupted packet passes with probability about $2^{-n}$, which is negligible for the CRC lengths used in practice (Section~\ref{sec:results} gives ours). The GS acknowledges only the packets it decodes, so an acknowledgment missing at the end of the slot tells the UAV that the packet was lost and the GS status is unchanged. We assume acknowledgments always arrive, since they carry only the status (and, after an upload, the ground label) and the GS sends them at higher power.

\begin{remark}[Undetected errors]\label{rem:flip}
Suppose instead that a decoded label is flipped with a small probability $\ee_{\rm u}$. Because the label is binary, a flipped disagreeing label equals the current GS status, so it changes nothing. A bit error on the link can therefore cancel an update but cannot make a correct GS status wrong. Its only effect is to replace $q_1(g)$ by $q_1(g)(1-\ee_{\rm u})$ in everything that follows. An undetected error in a raw frame corrupts the photo rather than the label, and its effect folds into the ground error $\ee_{\rm g}(c)$.
\end{remark}

\subsection{Energy Consumption}
A transmission costs the energy of the radio while it transmits plus a fixed overhead per packet. With the rateless code above, the radio is on for $\min\{L_a/C(h),\tau_{\max}\}$, so the expected energy of action $a$ is $P_{\rm r}\,\E_h\big[\min\{L_a/C(h),\tau_{\max}\}\big]+E_{\rm o}$. The controller needs this value before it transmits, when $h$ is unknown, and we approximate it by evaluating the airtime at the mean rate $\bar R(g)=W\log_2(1+\bar\gamma(g))$:
\begin{equation}\label{eq:energy}
E_a(g)=P_{\rm r}\min\Big\{\frac{L_a}{\bar R(g)},\tau_{\max}\Big\}+E_{\rm o},\quad a\in\{1,2\},
\end{equation}
and $E_0=0$. Here $P_{\rm r}$ is the power drawn by the radio and $E_{\rm o}$ is the overhead of waking up the radio, accessing the channel and receiving the acknowledgment. For a short label the overhead dominates and the approximation is accurate. For a raw frame it underestimates the cost most at intermediate distances, where a deep fade can extend the transmission to the full window; Section~\ref{sec:results} quantifies the gap in our scenario. Capture, inference and propulsion~\cite{zeng2019rotary} cost the same whatever the decision, so we constrain only the communication energy, through a long-run average budget of $\bar E$ joules per slot:
\begin{equation}\label{eq:budget}
\limsup_{T\to\infty}\frac1T\sum_{t=0}^{T-1}\E\Big[E_{a(t)}\big(g(t)\big)\Big]\le\bar E.
\end{equation}
Here $E_{a(t)}(g(t))$ is~\eqref{eq:energy} evaluated at the action and link state of slot $t$, which are random: they depend on the scene, the classifier outputs, the blockage and past deliveries. Up to the approximation in~\eqref{eq:energy}, \eqref{eq:budget} limits the long-run average energy that the radio draws. The $\limsup$ is needed because the running average need not converge for every policy. An average budget allows the UAV to spend more in some parts of the mission, for example near the GS where uploads are cheap, and less in others.

\subsection{The GS Status and the AoII}
Let $o(t)\in\{0,1\}$ indicate that the packet sent in slot $t$ is decoded, so $\Prob\{o(t)=1\mid a(t)=a,g(t)=g\}=q_a(g)$. The GS status, which we denote by $\hat X(t)$, is the most recent label the GS has received, whether it came from the UAV or from the ground classifier:
\begin{equation}\label{eq:est}
\hat X(t+1)=\begin{cases}\lab(t), & a(t)=1,\ o(t)=1,\\ \lab^{\rm g}(t), & a(t)=2,\ o(t)=1,\\ \hat X(t), & \text{otherwise.}\end{cases}
\end{equation}
Holding the last label is the natural estimator when the state is persistent ($p_x<\frac12$). Each acknowledgment carries the new status $\hat X(t+1)$, and no acknowledgment means that the status did not change, so the UAV always knows $\hat X(t)$. It never knows $X(t)$.

The GS status can be wrong for two reasons. The scene may change, for example when the UAV passes from flooded patches to dry ones, while the GS keeps the old status until the UAV updates it. Or the update itself may be wrong: the small on-board classifier mislabels some photos, especially ambiguous ones such as a dry field after rain, and a wrong label that reaches the GS turns a correct status into a wrong one. What matters to the response teams is how long the status has been wrong: a status that is wrong for one slot does little harm, while one that is wrong for a minute may send a team to the wrong place. We therefore measure how well the GS knows the scene by the AoII~\cite{maatouk2020aoii}. The AoII is zero while the GS status is correct and counts the slots for which it has been wrong otherwise. After the scene moves to $X(t+1)$, it evolves as
\begin{equation}\label{eq:aoii}
\Delta(t+1)=\big(\Delta(t)+1\big)\,\ind{X(t+1)\neq\hat X(t+1)},\quad \Delta(0)=0.
\end{equation}
The AoII captures two effects that the AoI misses. A delivered wrong label makes a correct GS wrong, so the AoII starts to grow right after a delivery. A GS that has been wrong for a long time is penalized more than one that has just become wrong, which favors correcting long errors first.

Table~\ref{tab:aoiiex} gives an example from the flood scenario over eight slots. In slot $2$ a wrong ``flooded'' label is delivered, so the GS becomes wrong in slot $3$ although the scene has not changed. In slot $4$ the patch floods, and the GS happens to become correct. In slot $4$ the UAV also delivers a wrong ``dry'' label, so the GS is wrong again in slots $5$ and $6$, until a correct label sent in slot $6$ fixes it.

\begin{table}[t]
\caption{AoII in a flood example ($1$: flooded, $0$: dry)}\label{tab:aoiiex}
\centering\footnotesize
\setlength{\tabcolsep}{5pt}
\begin{tabular}{@{}lcccccccc@{}}
\toprule
Slot $t$ & 1 & 2 & 3 & 4 & 5 & 6 & 7 & 8\\
\midrule
True state $X(t)$ & 0 & 0 & 0 & 1 & 1 & 1 & 1 & 1\\
Delivered label in slot $t$ & -- & 1 & -- & 0 & -- & 1 & -- & --\\
GS status $\hat X(t)$ & 0 & 0 & 1 & 1 & 0 & 0 & 1 & 1\\
AoII $\Delta(t)$ & 0 & 0 & 1 & 0 & 1 & 2 & 0 & 0\\
\bottomrule
\end{tabular}
\end{table}

The AoII is random, since the scene, the on-board errors and the deliveries are, and it depends on the true state, which the UAV never observes. Section~\ref{sec:problem} therefore takes its long-run average as the objective and replaces it with a belief that the UAV can compute.

\subsection{Order of Events}\label{sec:events}
Within slot $t$, events occur in the following order.
\begin{enumerate}
\item The UAV observes the link state $g(t)$, captures $Y(t)$ and computes $\lab(t)$ and $c(t)$.
\item It chooses $a(t)$ and transmits if $a(t)\neq0$.
\item The packet is decoded with probability $q_{a(t)}(g(t))$. If it is, the GS updates its status by~\eqref{eq:est} and sends the acknowledgment.
\item The scene moves to $X(t+1)$, and the AoII becomes $\Delta(t+1)$ by~\eqref{eq:aoii}.
\end{enumerate}
%% =====================================================================
\section{Problem Formulation}\label{sec:problem}
%% =====================================================================
This section states Q1 as a constrained optimization problem and reduces it to a belief about the AoII that the UAV can compute.

\subsection{Information and Policies}
When the UAV chooses $a(t)$, it knows its labels, confidences and link states up to slot $t$, and the outcome of every earlier transmission as reported by the acknowledgments,
\begin{equation}\label{eq:info}
\I(t)=\big\{\lab(s),c(s),g(s)\big\}_{s\le t}\cup\big\{a(s),\hat X(s+1)\big\}_{s<t},
\end{equation}
together with the ground label of every decoded upload.
It does not know the state $X(t)$ and therefore cannot compute the AoII $\Delta(t)$. A policy $\psi$ is a rule that maps the information $\I(t)$ to a probability distribution on the action set $\A$, from which $a(t)$ is drawn. We allow randomized policies because problems with an average constraint may need them to use the budget exactly. A transmission is pointless when the link is blocked, so we consider only policies with $a(t)=0$ whenever $q_1(g(t))=0$, and we denote their set by $\Psi$.

\subsection{The Optimization Problem}
The response teams need a GS status that is correct as often as possible and, when it is wrong, is corrected quickly. The long-run average AoII measures both: it grows with the fraction of slots in which the status is wrong, and it grows faster when the wrong periods are long, since the AoII counts the slots of each wrong period. Under a policy $\psi$, the average AoII and energy are
\begin{align}
\bar\Delta^{\psi}&=\limsup_{T\to\infty}\frac1T\sum_{t=0}^{T-1}\E^{\psi}\big[\Delta(t)\big],\label{eq:avgaoii}\\
\bar E^{\psi}&=\limsup_{T\to\infty}\frac1T\sum_{t=0}^{T-1}\E^{\psi}\Big[E_{a(t)}\big(g(t)\big)\Big],\label{eq:avgE}
\end{align}
where $\E^{\psi}$ is the expectation when the actions follow $\psi$, so that $\bar E^{\psi}$ is the left side of the budget~\eqref{eq:budget} under $\psi$. The UAV must keep the average AoII as low as possible without exceeding the budget:
\begin{equation}\label{eq:P}
\textbf{(P)}\qquad \min_{\psi\in\Psi}\ \bar\Delta^{\psi}\qquad\text{subject to}\qquad \bar E^{\psi}\le\bar E.
\end{equation}

Three features make (P) hard. First, it is partially observed: the cost $\Delta(t)$ depends on $X(t)$, which the UAV never sees, so the policy must act on its own inference about the state. Second, the cost has memory: by~\eqref{eq:aoii}, a wrong GS status adds to the AoII in every slot until it is corrected, so an action affects the cost of all the following slots, not only of the next one. Third, the energy constraint couples all slots: energy spent now is not available later, so the value of an update depends on the updates that the budget will allow in the future. We handle the first feature exactly, in the rest of this section, by replacing the unknown AoII with a probability distribution over its values. We handle the other two approximately: Section~\ref{sec:structure} values each action by its effect on the next slot, and Section~\ref{sec:lyap} prices energy by what has already been spent.

\subsection{A Belief on the AoII}\label{sec:belief}
The information $\I(t)$ grows with $t$. The following lemma shows that a much smaller object carries everything the UAV needs to predict the AoII.

\begin{lemma}[Sufficient statistic]\label{lem:suff}
Given $\hat X(t)$, the AoII determines the state: $X(t)=\hat X(t)$ if $\Delta(t)=0$ and $X(t)=1-\hat X(t)$ if $\Delta(t)>0$. Hence the pair $(\hat X(t),\pi_t)$, together with the link state $g(t)$ that the UAV observes, where
\begin{equation}\label{eq:belief}
\pi_t(\delta)=\Prob\big\{\Delta(t)=\delta\mid\I(t)\big\},\qquad\delta\in\{0,1,\dots\},
\end{equation}
summarizes $\I(t)$ for predicting the future AoII.
\end{lemma}
\begin{proof}
By~\eqref{eq:aoii}, $\Delta(t)>0$ iff $X(t)\neq\hat X(t)$, so $(\hat X(t),\Delta(t))$ determines the binary $X(t)$. By~\eqref{eq:source} and Assumption~\ref{as:inf}, the future scene and classifier outputs then depend on $\I(t)$ only through $\hat X(t)$ and $\pi_t$; the link evolves independently of the scene, and its current state $g(t)$ is observed.
\end{proof}

The mass of $\pi_t$ at $\delta=0$ is the probability that the GS is correct, and its mass at $\delta>0$ tells how long the GS has probably been wrong. For a given $\delta$, we write $x_\delta$ for the state implied by Lemma~\ref{lem:suff}, i.e., $x_0=\hat X(t)$ and $x_\delta=1-\hat X(t)$ for $\delta>0$.

To store the belief with finite memory, we truncate its support at $\Delta_{\max}$, adding all mass above $\Delta_{\max}$ to the mass at $\Delta_{\max}$. The belief is then a vector of $\Delta_{\max}+1$ numbers. The truncation is harmless as long as $\Delta_{\max}$ is well above the AoII values that actually occur, so that the mass at $\Delta_{\max}$ stays negligible; Section~\ref{sec:results} gives our value and checks this.

\subsection{Belief Update}\label{sec:update}
The belief changes for three reasons in each slot: the UAV sees a new on-board label, it may receive a ground label, and the scene evolves. Let $\pi^-_t$ denote the belief at the start of slot $t$, before the new frame is classified. The update has three steps.

\emph{Step 1: correction with the on-board label.} By Bayes' rule, the frame changes the belief in proportion to its likelihood under each value of the AoII:
\begin{equation}\label{eq:bayes}
\pi_t(\delta)\propto\pi^-_t(\delta)\,\Prob\big\{\lab(t),c(t)\mid\Delta(t)=\delta,\I(t-1)\big\}.
\end{equation}
Under Assumption~\ref{as:inf}(ii), the label and the confidence depend on the past only through the current state, which is $x_\delta$ by Lemma~\ref{lem:suff}. By~\eqref{eq:lik}, the likelihood in~\eqref{eq:bayes} is therefore proportional to $c_{x_\delta}(t)/\varpi_{x_\delta}$. All hypotheses $\delta>0$ imply the same state $1-\hat X(t)$, so the likelihood takes only two values, one if the GS is right ($\delta=0$) and one if it is wrong ($\delta>0$). Only their ratio matters. Let $d(t)=\ind{\lab(t)\neq\hat X(t)}$ indicate that the label disagrees with the GS. Their ratio is
\begin{equation}\label{eq:Lambda}
\Lambda_t=\frac{c_{1-\hat X(t)}(t)}{c_{\hat X(t)}(t)}\cdot\frac{\varpi_{\hat X(t)}}{\varpi_{1-\hat X(t)}},
\end{equation}
which says how many times more likely the frame is if the GS is wrong than if it is right. By~\eqref{eq:cx}, the first factor is $c(t)/(1-c(t))$ if the label disagrees with the GS and $(1-c(t))/c(t)$ if it agrees. A label that disagrees with the GS thus has $\Lambda_t$ increasing in its confidence: it is evidence that the GS is wrong, and the more so the more confident it is. A label that agrees with the GS is evidence that the GS is right.

Exact Bayes' rule would multiply the mass at $\delta>0$ by $\Lambda_t$ relative to the mass at $\delta=0$. This treats every frame as an independent observation, which overstates the evidence when errors come in runs. If a frame overlaps the previous one and the classifier missed the flood in the previous frame, a second miss is far from independent confirmation (Fig.~\ref{fig:inference}(c)). We therefore temper the ratio with an exponent $\nu\in(0,1]$. This is a practical correction, not a consequence of the model, and $\nu$ is chosen on data. With $\mu^-_t=\pi^-_t(\Delta>0)$ the prior probability that the GS is wrong, the corrected belief is
\begin{equation}\label{eq:post}
\pi_t(\delta)=\frac{\pi^-_t(\delta)}{1-\mu^-_t+\mu^-_t\Lambda_t^{\nu}}\times\begin{cases}1, & \delta=0,\\ \Lambda_t^{\nu}, & \delta>0.\end{cases}
\end{equation}
The exponent counts each frame as a fraction $\nu$ of an independent observation, and $\nu=1$ recovers Bayes' rule. A rough guide is the inverse of the mean length of an error run. A guide with a firmer basis comes from the correlation of the errors. Suppose the error indicators of consecutive frames form a stationary sequence whose correlation at lag $k$ is $\kappa^k$. The sum of $n$ such indicators then has about $(1+\kappa)/(1-\kappa)$ times the variance of a sum of $n$ independent ones, so $n$ frames carry roughly the evidence of $n(1-\kappa)/(1+\kappa)$ independent frames, which suggests $\nu=(1-\kappa)/(1+\kappa)$. In the evaluation we select $\nu$ on tuning data together with the other parameters of the controller, and Section~\ref{res:robust} reports the results with the correlation rule instead. We also include, as an option, the symmetric update that ignores the prior and uses $\varpi_0=\varpi_1$ with $\nu=1$, in which the confidence itself is the likelihood.

Step 1 multiplies the odds that the GS is wrong by $\Lambda_t^{\nu}$. It moves mass between ``the GS is right'' and ``the GS is wrong'' but does not change the shape of the belief within $\delta>0$.

\emph{Step 2: correction with the ground label.} If the UAV uploaded the frame and the GS decoded it, the acknowledgment carries the ground label, and~\eqref{eq:post} is applied once more with $\nu=1$. Its error probability $\ee_{\rm g}=\ee_{\rm g}(c(t))$ is defined given the state, so no prior correction is needed, and the ratio is $(1-\ee_{\rm g})/\ee_{\rm g}$ if the ground label disagrees with $\hat X(t)$ and $\ee_{\rm g}/(1-\ee_{\rm g})$ otherwise.

\emph{Step 3: prediction.} After the transmission, the GS holds $\xh'=\hat X(t+1)$, which the UAV knows from the acknowledgment or from the absence of one. For each $\delta$, the state $X(t)=x_\delta$ is known, so the UAV knows whether the GS status is now correct for that hypothesis. By~\eqref{eq:aoii}, the AoII grows to $\delta+1$ if $X(t+1)\neq\xh'$ and resets to 0 otherwise. If $x_\delta=\xh'$, the status becomes wrong only if the state changes, which happens with probability $p_{x_\delta}$. If $x_\delta\neq\xh'$, it stays wrong if the state does not change, with probability $1-p_{x_\delta}$. Hence
\begin{equation}\label{eq:pred}
\Prob\big\{\Delta(t+1)=\delta+1\mid\Delta(t)=\delta\big\}=\begin{cases}p_{x_\delta}, & x_\delta=\xh',\\ 1-p_{x_\delta}, & x_\delta\neq\xh',\end{cases}
\end{equation}
with $\Delta(t+1)=0$ otherwise. Summing over $\delta$ under $\pi_t$ gives $\pi^-_{t+1}$. A delivered label that flips the GS status swaps the two cases of~\eqref{eq:pred}, which is how a correction moves the mass of the belief back towards $\delta=0$.

Two summaries of the belief appear in the rest of the paper:
\begin{equation}\label{eq:musum}
\mu_t=\pi_t(\Delta>0),\qquad S_t=\E_{\pi_t}\big[(\Delta+1)\ind{\Delta>0}\big].
\end{equation}
The first is the probability that the GS is wrong. The second weighs each wrong hypothesis by the AoII it will have in the next slot if the GS stays wrong. We write $\mu^-_t$ and $S^-_t$ for the same summaries of $\pi^-_t$.

\subsection{A Worked Example}
\begin{example}\label{ex:belief}
Let the GS report ``not flooded'' ($\hat X(t)=0$), and let the UAV hold the belief $\pi^-_t(0)=0.95$, $\pi^-_t(1)=0.05$: it is fairly sure the GS is right. A frame then arrives with label ``flooded'' and confidence $c=0.9$, so $d(t)=1$. With $\varpi_0=0.81$ and $\varpi_1=0.19$,~\eqref{eq:Lambda} gives $\Lambda_t=(0.81/0.19)(0.9/0.1)=38.4$: the frame is $38$ times more likely if the GS is wrong. With exact Bayes' rule ($\nu=1$),~\eqref{eq:post} gives $\mu_t=0.05\cdot38.4/(0.95+0.05\cdot38.4)=0.67$: one confident frame makes the GS more likely wrong than right. With $\nu=0.5$ the same frame gives only $\mu_t=0.25$, and with $\nu=0.25$ it gives $\mu_t=0.12$. Suppose the UAV stays silent. The prediction~\eqref{eq:pred} with $p_0=0.039$ and $p_1=0.166$ gives $\mu^-_{t+1}=0.23$ for $\nu=0.5$. If the next frame again says ``flooded'' with $c=0.9$, the belief rises to $\mu_{t+1}=0.65$. With the tempered update, the UAV therefore needs two consistent frames to become as convinced as one frame makes it under exact Bayes' rule. Section~\ref{sec:structure} continues this example to show what the UAV does in each of these two slots.
\end{example}

\subsection{Equivalent Belief-State Problem and Solution Approach}
By Lemma~\ref{lem:suff}, (P) is equivalent to a constrained Markov decision process whose state is the belief together with the observable quantities, $(\pi^-_t,\hat X(t),\lab(t),c(t),g(t))$. The per-slot cost is the expected AoII $\E[\Delta(t+1)\mid\I(t),a(t)]$, which the belief gives, and the per-slot energy is $E_{a(t)}(g(t))$. The transition kernel follows from the model: the next label and confidence are drawn from their law given $X(t+1)$; $g(t+1)$ follows the trajectory and the blockage chain~\eqref{eq:block}; $\hat X(t+1)$ follows~\eqref{eq:est}; and $\pi^-_{t+1}$ follows the three steps of Section~\ref{sec:update}.

The standard approach prices the energy constraint with a Lagrange multiplier $\lambda\ge0$, solves the unconstrained problem by value iteration, and searches over $\lambda$ to meet the budget. This is impractical on a UAV: the belief lives in a high-dimensional simplex, and the kernel and the right multiplier depend on the law of the confidence, the blockage and the trajectory, which change between missions. Instead, Section~\ref{sec:structure} characterizes the value of each action in one slot in closed form, and Section~\ref{sec:lyap} builds an online policy on this characterization, with an energy price that adapts to the energy spent.
%% =====================================================================
\section{The Value of Each Action}\label{sec:structure}
%% =====================================================================
Any reasonable policy must compare, in each slot, how much each action is expected to lower the AoII with how much energy it costs. This section derives the first quantity in closed form and shows what it implies for the decision. We value an action by its effect on the AoII of the next slot, because that effect depends only on the current belief; the effect on later slots depends on future decisions as well, and Section~\ref{sec:lyap} leaves it to the energy price. Throughout, we fix slot $t$, let $\xh=\hat X(t)$, and write
\begin{equation}\label{eq:rw}
r=p_{\xh},\qquad w=p_{1-\xh}.
\end{equation}
In words, $r$ is the probability that the state leaves the value the GS currently reports, and $w$ is the probability that it leaves the other value. By~\eqref{eq:source}, both are below $\frac12$.

\subsection{One-Step Expected AoII}
We first compute the expected AoII in slot $t+1$ for a given hypothesis $\Delta(t)=\delta$ and a given new status $\xh'=\hat X(t+1)$. By~\eqref{eq:aoii}, $\Delta(t+1)$ is either $\delta+1$ or $0$, so its expectation is $\delta+1$ times the probability in~\eqref{eq:pred}: $(\delta+1)p_{x_\delta}$ if the new status equals $x_\delta$, and $(\delta+1)(1-p_{x_\delta})$ otherwise. The three actions differ only in which status $\xh'$ they produce, and with what probability. We evaluate this expectation for each of them, assuming the packet is delivered, and average it over $\pi_t$.

\emph{Silence.} The status stays $\xh'=\xh$. If $\delta=0$, then $x_0=\xh=\xh'$ and the probability is $p_{\xh}=r$. If $\delta>0$, then $x_\delta=1-\xh\neq\xh'$ and the probability is $1-p_{1-\xh}=1-w$. Averaging over $\pi_t$ and using~\eqref{eq:musum},
\begin{equation}\label{eq:J0}
J_0=(1-\mu_t)\,r+S_t\,(1-w).
\end{equation}

\emph{Label.} If $d(t)=0$, the label equals $\xh$ and changes nothing, so the result is $J_0$. If $d(t)=1$, a delivered label flips the GS status to $\xh'=1-\xh$. Now the case $\delta=0$ has $x_0\neq\xh'$ and probability $1-r$, and the case $\delta>0$ has $x_\delta=\xh'$ and probability $w$, so
\begin{equation}\label{eq:J1}
J_1=(1-\mu_t)(1-r)+S_t\,w.
\end{equation}

\emph{Raw upload.} The ground label equals $X(t)$ with probability $1-\ee_{\rm g}$, where $\ee_{\rm g}=\ee_{\rm g}(c(t))$, and by Assumption~\ref{as:inf}(iii) this event is independent of $\Delta(t)$ given $c(t)$. If the ground label is correct, then $\xh'=x_\delta$ and the probability is $p_{x_\delta}$, which is $r$ for $\delta=0$ and $w$ for $\delta>0$. If it is wrong, the probability is $1-p_{x_\delta}$. Mixing the two cases,
\begin{equation}\label{eq:J2}
J_2=(1-\mu_t)\big[r+\ee_{\rm g}(1-2r)\big]+S_t\big[w+\ee_{\rm g}(1-2w)\big].
\end{equation}

A transmission is delivered with probability $q_a(g)$ and has no effect otherwise, and delivery is independent of the scene given $g$. The expected AoII after action $a$ is therefore $(1-q_a)J_0+q_aJ_a=J_0-q_a(J_0-J_a)$. We call the difference $G_a=J_0-J_a$ the gain of action $a$: it is the reduction in the expected AoII of the next slot that a delivered packet of type $a$ brings. Subtracting~\eqref{eq:J1} and~\eqref{eq:J2} from~\eqref{eq:J0} proves the following lemma.

\begin{lemma}[One-step expected AoII]\label{lem:onestep}
Under Assumption~\ref{as:inf}, $\E[\Delta(t+1)\mid\I(t),a(t)=a]=J_0-q_a(g)\,G_a$, with $G_0=0$ and
\begin{align}
G_1&=\big[S_t(1-2w)-(1-\mu_t)(1-2r)\big]\,d(t), \label{eq:G1}\\
G_2&=(1-\ee_{\rm g})\,S_t(1-2w)-\ee_{\rm g}\,(1-\mu_t)(1-2r). \label{eq:G2}
\end{align}
\end{lemma}

Lemma~\ref{lem:onestep} is exact when $\pi_t$ is the posterior, that is, with $\nu=1$ and no truncation. With $\nu<1$ the controller evaluates the same expressions at the tempered belief~\eqref{eq:post}, so $J_0-q_aG_a$ is then its estimate of the expected AoII; the structural results below hold for any $\nu\in(0,1]$.

Both gains are built from the same two terms. The term $S_t(1-2w)$ is the AoII that an update removes when the GS is wrong and the update is right. The term $(1-\mu_t)(1-2r)$ is the AoII that an update creates when the GS is right and the update is wrong. Since $r,w<\frac12$, both terms are nonnegative. A delivered disagreeing label removes the first term in full, because a label that disagrees with a wrong GS is necessarily correct. It also creates the second in full, because a label that disagrees with a correct GS is necessarily wrong. A ground label removes the first term only with probability $1-\ee_{\rm g}$, but it creates the second only with probability $\ee_{\rm g}$. The label is therefore a bold update and the raw upload a cautious one. A label that is probably wrong has $G_1<0$ and would increase the expected AoII, while the raw upload has $G_2>0$ as soon as $S_t(1-2w)(1-\ee_{\rm g})>\ee_{\rm g}(1-\mu_t)(1-2r)$, which holds even for a small probability that the GS is wrong when the ground classifier is accurate.

\setcounter{example}{0}
\begin{example}[continued]
In slot $t$ of Example~\ref{ex:belief} with $\nu=0.5$, the belief is $\mu_t=0.25$ and $S_t=0.49$, and the ground error at $c=0.9$ is $\ee_{\rm g}=0.028$ (Fig.~\ref{fig:inference}(b)). With $\xh=0$, $r=p_0=0.039$ and $w=p_1=0.166$. Lemma~\ref{lem:onestep} gives $G_1=-0.37$ and $G_2=0.30$: sending the label would more likely harm than help, while asking the ground would help. In slot $t+1$, after the second disagreeing frame, $\mu_{t+1}=0.65$ and the gains become $G_1=0.94$ and $G_2=1.21$. Now the GS is probably wrong, and the label is worth sending.
\end{example}

For a symmetric source the gains take a simpler form, in terms of the mean of the belief.

\begin{corollary}[Symmetric source]\label{cor:sym}
If $p_0=p_1=p$, then with $\Dbar_t=\E_{\pi_t}[\Delta(t)]$,
\begin{align*}
G_1&=(1-2p)\big(\Dbar_t+2\mu_t-1\big)\,d(t),\\
G_2&=(1-2p)\big(\mu_t+(1-\ee_{\rm g})\Dbar_t-\ee_{\rm g}\big).
\end{align*}
\end{corollary}
\begin{proof}
With $r=w=p$, both terms of the gains carry the factor $1-2p$. Since $S_t=\E_{\pi_t}[\Delta\ind{\Delta>0}]+\mu_t=\Dbar_t+\mu_t$, substituting in~\eqref{eq:G1} and~\eqref{eq:G2} gives the result.
\end{proof}

Corollary~\ref{cor:sym} gives a simple test for the label: a disagreeing label lowers the expected AoII if and only if $\Dbar_t+2\mu_t>1$. When the GS has just become wrong, $\Dbar_t\approx\mu_t$, and the test reduces to $\mu_t>\frac13$. As the GS stays wrong, $\Dbar_t$ grows and a smaller probability of error suffices. The test does not depend on $p$, since the factor $1-2p$ only scales both gains.

\subsection{Agreeing Labels}
The gains also show which updates are never useful.

\begin{proposition}[No confirmations]\label{prop:confirm}
If $d(t)=0$, then $G_1=0$. Hence a policy that minimizes the expected AoII plus a nonnegative energy price never needs to send a label that agrees with the GS.
\end{proposition}
\begin{proof}
If $d(t)=0$, a delivered label leaves the status unchanged, so $J_1=J_0$ and $G_1=0$. Sending it then costs $E_1(g)>0$ and changes nothing else, so silence does at least as well for any energy price $\lambda\ge0$.
\end{proof}

Proposition~\ref{prop:confirm} rules out a common design in which the device reports every label so that the receiver can confirm its status; the pre-identification baseline of Section~\ref{sec:results} spends most of its energy on such labels.

\subsection{Threshold Structure of the Per-Slot Decision}
Lemma~\ref{lem:onestep} measures the value of an action in AoII, while its cost is measured in energy. A per-slot rule must convert one into the other. We consider rules of the form
\begin{equation}\label{eq:rule}
a(t)\in\arg\max_{a\in\A}\big\{V\,q_a(g)\,G_a-\lambda(t)\,E_a(g)\big\},
\end{equation}
where $V>0$ weighs the AoII and $\lambda(t)\ge0$ is a price per joule that may change over time. Section~\ref{sec:lyap} shows that the drift-plus-penalty method produces exactly this rule, with $\lambda(t)$ equal to an energy debt. Here we ask how the decision depends on the confidence of the frame.

The confidence enters the gains only through the belief $\pi_t$, and by~\eqref{eq:post} only through the ratio $\Lambda_t$. By~\eqref{eq:Lambda}, $\Lambda_t$ increases in $c$ if $d(t)=1$ and decreases in $c$ if $d(t)=0$. The next lemma shows that both gains, and any nonnegative combination of their two terms, increase with $\Lambda_t$.

\begin{lemma}[Monotonicity in the evidence]\label{lem:mono}
Fix $\pi^-_t$. For any $A,B\ge0$, the quantity $AS_t-B(1-\mu_t)$ is a nondecreasing function of $\Lambda_t$.
\end{lemma}

The proof is in Appendix~\ref{app:mono}. Combined with Lemma~\ref{lem:onestep}, it yields the structure of the decision.

\begin{proposition}[Confidence thresholds]\label{prop:threshold}
Fix $\pi^-_t$, $\hat X(t)$, $g$ and $\lambda(t)$, and let $I$ be an interval of confidences on which $\ee_{\rm g}(c)$ is constant, such as one bin of the estimate used in Section~\ref{sec:results}. For $c\in I$:
\begin{enumerate}
\item[(a)] If $d(t)=1$, there are thresholds $c_{\rm r}\le c_{\lab}$ such that~\eqref{eq:rule} stays silent for $c<c_{\rm r}$, uploads the frame for $c_{\rm r}\le c<c_{\lab}$, and sends the label for $c\ge c_{\lab}$. Any of the three intervals may be empty.
\item[(b)] If $d(t)=0$, it uploads the frame if and only if $c\le c^0_{\rm r}$ for some threshold $c^0_{\rm r}$, and stays silent otherwise.
\item[(c)] For fixed $\mu^-_t$, the silent regions shrink as $S^-_t$ grows, and they grow with $\lambda(t)$.
\end{enumerate}
\end{proposition}

The proof is in Appendix~\ref{app:threshold}. Across bins the thresholds can shift where $\ee_{\rm g}(c)$ jumps, which causes the steps in Fig.~\ref{fig:regions}. Part (a) is the rule of the title. A frame that disagrees with the GS is sent as a label when the UAV is sure of it, sent to the ground when the UAV is unsure, and withheld when it is too unsure for either to pay off. Part (b) says that a frame that agrees with the GS can still be worth uploading, but only if it is unsure: an unsure agreement leaves room for the GS to be wrong, and the ground classifier can check. Part (c) says that the UAV acts more readily the longer the GS has probably been wrong, and less readily as energy becomes costlier.

\setcounter{example}{0}
\begin{example}[continued]
Let the UAV be $1.2$\,km from the GS, where $q_1\approx1$, $q_2=0.78$, $E_1=20$\,mJ and $E_2=190$\,mJ. In slot $t$, the label has $q_1G_1<0$, and the upload has $q_2G_2=0.23$, so~\eqref{eq:rule} uploads the frame if $\lambda/V<0.23/0.19=1.2$\,J$^{-1}$ and stays silent otherwise. In slot $t+1$, $q_1G_1=0.94$ and $q_2G_2=0.95$ are almost equal, but the label costs a tenth of the upload, so the UAV sends the label for any price between about $0.05$ and $47$\,J$^{-1}$. At the price $\lambda/V=2$\,J$^{-1}$ used in Fig.~\ref{fig:regions}, the UAV thus waits in slot $t$ and sends the label in slot $t+1$.
\end{example}

\begin{figure}[t]
\centering
\includegraphics[width=0.8\columnwidth]{figures/regions.pdf}
\caption{Example of the decisions of~\eqref{eq:rule} versus the slots since the last delivery and the on-board confidence $c$. The GS reports ``not flooded'' with $5\%$ initial doubt, the UAV is $1.2$\,km from the GS ($q_1\approx1$, $q_2=0.78$, $E_1=20$\,mJ, $E_2=190$\,mJ), $\lambda/V=2$\,J$^{-1}$ and $\nu=0.25$. (a) An agreeing label is never sent; the frame is uploaded once the status is old, unless the UAV is very confident. (b) A disagreeing frame is uploaded when the UAV is unsure and sent as a label when it is confident (Proposition~\ref{prop:threshold}).}
\label{fig:regions}
\end{figure}

Fig.~\ref{fig:regions} shows the decision regions of~\eqref{eq:rule} for the scenario of the example, as a function of the slots since the last delivery and of the confidence. It uses the confidence-dependent $\ee_{\rm g}(c)$ estimated from our data, which Proposition~\ref{prop:threshold} does not cover; the order silent, raw, label still holds. When the label disagrees with the GS, the confidence needed to send it falls as the GS status ages, because the prior probability that the GS is wrong grows. When the label agrees with the GS but is unsure, the UAV uploads the frame once the status is old enough, which amounts to asking the GS to verify its status.
%% =====================================================================
\section{Online Policy: CA-DPP}\label{sec:lyap}
%% =====================================================================
The rule~\eqref{eq:rule} needs an energy price $\lambda(t)$. A fixed price would have to be tuned to the mission, and a wrong value either wastes the budget or leaves it unused. This section derives a price from the energy actually spent, using Lyapunov drift-plus-penalty~\cite{neely2010book}, and shows that the resulting policy meets the budget.

\subsection{Energy Debt as a Virtual Queue}
We track how far the UAV has spent beyond its budget with a virtual queue,
\begin{equation}\label{eq:Z}
Z(t+1)=\max\Big\{Z(t)+E_{a(t)}\big(g(t)\big)-\bar E,\ 0\Big\},\qquad Z(0)=0.
\end{equation}
In each slot the queue grows by the energy spent and drains by the budget, so $Z(t)$ is the energy debt of the UAV. Since $\max\{z,0\}\ge z$, summing~\eqref{eq:Z} over $t<T$ gives $Z(T)\ge\sum_{t<T}\big(E_{a(t)}(g(t))-\bar E\big)$, or
\begin{equation}\label{eq:Zbound}
\frac1T\sum_{t=0}^{T-1}E_{a(t)}\big(g(t)\big)\le\bar E+\frac{Z(T)}{T}.
\end{equation}
Hence the budget~\eqref{eq:budget} holds whenever the debt grows slower than linearly. The policy should therefore keep $Z(t)$ small while it keeps the AoII low.

\subsection{Drift-Plus-Penalty}
To keep the queue small, we measure its size by $L(t)=\frac12Z(t)^2$ and its expected growth in one slot by the drift $\E[L(t+1)-L(t)\mid\I(t)]$. Squaring~\eqref{eq:Z} and using $\max\{z,0\}^2\le z^2$,
\[
Z(t+1)^2\le Z(t)^2+\big(E_{a(t)}-\bar E\big)^2+2Z(t)\big(E_{a(t)}-\bar E\big).
\]
The middle term is at most $\max\{E_{\max},\bar E\}^2$, where $E_{\max}=\max_{a,g}E_a(g)$. Taking conditional expectations,
\begin{equation}\label{eq:drift}
\E\big[L(t+1)-L(t)\,|\,\I(t)\big]\le C+Z(t)\,\E\Big[E_{a(t)}\big(g(t)\big)-\bar E\,\big|\,\I(t)\Big],
\end{equation}
with $C=\frac12\max\{E_{\max},\bar E\}^2$. The drift-plus-penalty method chooses, in each slot, the action that minimizes the right-hand side of~\eqref{eq:drift} plus $V$ times the expected penalty, here the expected AoII of the next slot. By Lemma~\ref{lem:onestep}, this is
\[
\min_{a\in\A}\Big\{Z(t)E_a(g)+V\big(J_0-q_a(g)G_a\big)\Big\}.
\]
The terms $C$, $Z(t)\bar E$ and $VJ_0$ do not depend on the action, so the minimization is equivalent to~\eqref{eq:rule} with $\lambda(t)=Z(t)$. The energy price is thus the energy debt. When the UAV has spent more than its budget, energy becomes expensive and the silent regions of Proposition~\ref{prop:threshold}(c) grow. When it has spent less, the debt drains to zero and every useful update is sent.

We call the resulting policy confidence-aware drift-plus-penalty (CA-DPP). Algorithm~\ref{alg:dpp} lists its steps.

\begin{algorithm}[t]
\caption{CA-DPP, executed by the UAV in slot $t$}\label{alg:dpp}
\begin{algorithmic}[1]
\State Classify $Y(t)$ to get $\lab(t)$, $c(t)$; observe $g(t)$; set $d(t)$
\State Correct $\pi^-_t\to\pi_t$ by~\eqref{eq:post}; compute $\mu_t$, $S_t$
\State Compute $G_1$, $G_2$ by~\eqref{eq:G1}--\eqref{eq:G2}
\If{$q_1(g(t))=0$} $a(t)\gets0$
\Else\ $a(t)\gets\arg\max_a\{Vq_a(g)G_a-Z(t)E_a(g)\}$, ties to silence, then to upload
\EndIf
\State Transmit; read the acknowledgment: $\hat X(t+1)$, and $\lab^{\rm g}(t)$ after a decoded raw upload
\State If a ground label arrived, correct $\pi_t$ with it (Step 2)
\State Predict $\pi^-_{t+1}$ by~\eqref{eq:pred}
\State $Z(t+1)\gets\max\big\{Z(t)+E_{a(t)}(g(t))-\bar E,0\big\}$
\end{algorithmic}
\end{algorithm}

\subsection{Energy Guarantee}
The next theorem shows that CA-DPP keeps the energy debt bounded on every sample path, so it meets the budget without any knowledge of the statistics of the scene, the classifier or the link.

\begin{theorem}[Budget]\label{thm:budget}
Let $E_{\min}=\min_{a\neq0,g}E_a(g)>0$. Under CA-DPP, for every sample path and every $t$,
\begin{equation}\label{eq:Zmax}
Z(t)\le Z_{\max}=\frac{V(\Delta_{\max}+1)}{E_{\min}}+E_{\max}.
\end{equation}
Consequently, $\frac1T\sum_{t<T}E_{a(t)}(g(t))\le\bar E+Z_{\max}/T$ for every $T$, and~\eqref{eq:budget} holds.
\end{theorem}

The proof is in Appendix~\ref{app:budget}. A gain never exceeds $\Delta_{\max}+1$, so once $Z(t)E_a(g)$ exceeds $V(\Delta_{\max}+1)$ for every transmission, the UAV stays silent and the debt drains. The bound~\eqref{eq:Zmax} is loose; it shows that the budget is met by construction, and Section~\ref{sec:results} shows how close the spent energy stays to the budget in practice.

Theorem~\ref{thm:budget} guarantees the energy on every sample path. For the AoII, CA-DPP is a heuristic with a clear basis: it minimizes the next-slot AoII plus the priced energy. We have no analytical bound on its gap to the optimum of (P) (Remark~\ref{rem:opt}), so Section~\ref{sec:results} assesses its AoII on data and bounds the gap numerically under Assumption~\ref{as:inf}.

\subsection{Complexity and Parameters}
A slot of CA-DPP costs $O(\Delta_{\max})$ operations, linear in the length of the belief vector. The memory is the belief vector, $\Delta_{\max}+1$ numbers, plus small tables of $\ee_{\rm g}(c)$, $q_a(g)$ and $E_a(g)$.

CA-DPP has two parameters, chosen on tuning data before the mission together with the statistics $p_0$, $p_1$, $\varpi$ and $\ee_{\rm g}(c)$. The weight $V$ trades the AoII against the energy debt: a larger $V$ lowers the AoII but lets the debt grow larger, so a finite mission overspends more. The exponent $\nu$ sets how much each frame moves the belief.

\subsection{Extensions and Performance}
\begin{remark}[Harvesting battery]\label{rem:harvest}
If the UAV, or a ground sensor running the same policy, harvests energy into a finite battery with content $B(t)$, the constraint becomes that the battery never runs empty. The perturbed Lyapunov method of~\cite{huang2013eh} applies the same derivation to the battery, shifted by a constant $\theta$, and yields~\eqref{eq:rule} with $\lambda(t)=\theta-B(t)$: energy becomes expensive when the battery is low. With a battery capacity of at most $\theta$, the price stays nonnegative, so Proposition~\ref{prop:confirm} still applies. By the argument behind Theorem~\ref{thm:budget}, choosing $\theta\ge E_{\max}+V(\Delta_{\max}+1)/E_{\min}$ ensures that the UAV never chooses an action that the battery cannot pay for.
\end{remark}

\begin{remark}[Optimality]\label{rem:opt}
For problems in which the per-slot cost depends only on the current state and action, drift-plus-penalty comes within $O(1/V)$ of the optimal cost~\cite{neely2010book}. This bound does not apply here, because the AoII depends on past decisions through the GS status. A version that applies~\eqref{eq:rule} between deliveries fits the renewal framework of~\cite{neely2013renewal}, in which each delivery starts a new frame; its bound is left to future work. Section~\ref{res:opt} bounds the gap numerically.
\end{remark}
%% =====================================================================
\section{Trace-Driven Evaluation}\label{sec:results}
%% =====================================================================
We evaluate CA-DPP on real classifier outputs, replayed over a simulated air-to-ground link. Replaying real outputs matters here: the on-board errors in the traces are correlated in time and only approximately calibrated, so they test the policy under conditions that violate Assumption~\ref{as:inf}. Section~\ref{res:budget} compares CA-DPP with the baselines, Sections~\ref{res:gains} and~\ref{res:sens} answer Q2 and Q3, and Section~\ref{res:V} examines the weight $V$ and the trajectory. Section~\ref{res:opt} bounds the distance to the optimal policy, Section~\ref{res:corr} isolates the effect of correlated errors, and Section~\ref{res:robust} collects robustness checks.

\subsection{Data}
We build the traces from the hand-labeled part of Sen1Floods11~\cite{bonafilia2020sen1floods11}. It contains 446 Sentinel-2 chips of $512\times512$ pixels at $10$\,m resolution from eleven flood events on several continents, with per-pixel water labels and official train, validation and test splits. No UAV flood dataset with pixel labels of comparable size exists, so this imagery stands in for the UAV view.

To turn a chip into a sequence of frames, a $64\times64$-pixel footprint scans the chip along a serpentine path with a stride of $32$ pixels, one footprint per slot. Consecutive footprints therefore overlap by half, as consecutive frames of a UAV flying at constant speed would. This scan gives the sequence of photos; the position of the UAV, which sets the link, follows the route described below and is generated separately (see Section~\ref{sec:discussion}). A footprint is labeled flooded if at least $10\%$ of its labeled pixels are water, and footprints with less than $80\%$ labeled pixels are skipped. This leaves $46{,}774$, $16{,}612$ and $16{,}811$ slots in the train, validation and test splits. The resulting state sequence is strongly correlated: on the training split, $p_0=0.039$ and $p_1=0.166$, so dry stretches last $26$ slots on average, flooded ones $6$, and $\varpi_1=19\%$ of footprints are flooded.

\subsection{Classifiers}\label{res:cls}
The on-board model must fit a small processor. It is a logistic regression on the blue, green, red and near-infrared bands averaged over $8\times8$ pixel blocks ($257$ parameters). The larger ground model is a gradient-boosted tree ensemble with $300$ trees; its $77$ inputs are the mean, standard deviation and $10/50/90$th percentiles of all $13$ bands and two water indices, plus the share of water pixels by each index. Both models are trained on the training chips and calibrated by isotonic regression on the validation chips, never on the test chips.

For weaker on-board models, we add Gaussian noise of $2$, $4$, $6$ and $8$ times each band's standard deviation to the inputs, then retrain and recalibrate. On the $174$ validation and test chips ($33{,}423$ slots, $20.5\%$ flooded), the base model and the four noisy variants are correct on $91.8\%$, $89.4\%$, $86.7\%$, $83.9\%$ and $81.8\%$ of the slots, and the ground model on $97.6\%$. Noise makes the errors lopsided: the noisiest variant labels $7\%$ of the test slots as flooded, against $21\%$ in truth, so $88\%$ of its errors are missed floods.

On the test chips, the expected calibration error (ECE, ten equal-width bins) is $1.4\%$ for the base on-board model, $2.2\%$ for the noisiest variant, which therefore still reports honestly when it is unsure (Fig.~\ref{fig:inference}(a)), and $0.2\%$ for the ground model. The on-board errors come in runs: over the validation and test chips the base model is wrong in $8.2\%$ of the slots but in $56\%$ of the slots that follow a wrong label, and a run lasts $2.2$ slots on average. This suggests $\nu\approx0.45$ by the run length; the lag-one correlation of the error indicator is $\kappa\approx0.5$, so the correlation rule of Section~\ref{sec:update} gives $\nu$ between $0.31$ and $0.35$ on the tuning folds. The tuning picks $\nu\in\{0.25,0.5\}$. The ground model is wrong on $20\%$ of the frames that the on-board model gets wrong and on $0.7\%$ of the others, and its error is about $8\%$ for $c<0.8$ and below $1\%$ for $c>0.95$.

\subsection{Scenario}
Table~\ref{tab:params} lists the scenario parameters. Over one orbit of $200$\,s, the distance to the GS ranges from $0.2$ to $2.8$\,km. The link parameters give a label energy of about $20$\,mJ everywhere and a raw upload that costs $58$ to $420$\,mJ and succeeds with probability $0.99$ to $0.25$, depending on the distance. Averaging over the fading instead of using~\eqref{eq:energy} gives $20$ to $22$\,mJ for a label and $64$ to $395$\,mJ for a raw frame, so~\eqref{eq:energy} is within $-30\%$ to $+12\%$ of the exact cost of a raw frame. Every packet carries a $32$-bit CRC. Blockage has $\rho_{\rm b}=0.02$ and $\rho_{\rm u}=0.1$, so $17\%$ of the slots are blocked, for $10$ slots on average. We consider budgets $\bar E$ from $0.25$ to $10$\,mJ per slot. At the lowest budget the UAV can afford a label in about one slot in $80$, and at the highest it can afford a raw upload in most slots near the GS. We truncate the belief at $\Delta_{\max}=300$, i.e., $1.2$\,kB in single precision. A slot then takes about two thousand arithmetic operations, one pass over the belief for each of the correction, the two gains and the prediction; on a microcontroller with a floating-point unit at $100$\,MHz this is on the order of tens of microseconds (our estimate at one operation per cycle), against a slot of $0.5$\,s. The value only needs to exceed the AoII values that occur: the longest AoII under CA-DPP in our runs is below $40$ slots, and runs with $\Delta_{\max}$ from $50$ to $1000$ give identical decisions.

\begin{table}[t]
\caption{Scenario parameters}\label{tab:params}
\centering\footnotesize
\setlength{\tabcolsep}{4pt}
\begin{tabular}{@{}ll@{}}
\toprule
Parameter & Value\\
\midrule
Slot length $\tau$; transmission window $\tau_{\max}$ & $0.5$\,s; $0.4$\,s\\
UAV altitude; orbit radius; center to GS & $100$\,m; $1.3$\,km; $1.5$\,km\\
Orbit period & $400$ slots\\
Carrier; bandwidth $W$; transmit power & $2.4$\,GHz; $1$\,MHz; $20$\,dBm\\
Noise figure; environment & $7$\,dB; urban~\cite{alhourani2014lap}\\
Blockage $\rho_{\rm b}$; $\rho_{\rm u}$ & $0.02$; $0.1$\\
Label size $L_1$; frame size $L_2$ & $32$\,B; $50$\,kB\\
Radio power $P_{\rm r}$; overhead $E_{\rm o}$ & $1$\,W; $20$\,mJ\\
Source $p_0$; $p_1$ (training split) & $0.039$; $0.166$\\
Belief support $\Delta_{\max}$ & $300$\\
Energy budget $\bar E$ & $0.25$ to $10$\,mJ/slot\\
\bottomrule
\end{tabular}
\end{table}

\begin{table*}[t]
\caption{Trace-driven results at two energy budgets (mean $\pm$ 95\% confidence half-width over 20 held-out runs: 5-fold cross-validation over chips, 4 runs per fold). $\Delta$: change in average AoII obtained by CA-DPP relative to the policy in that row (negative means CA-DPP is lower), with paired 95\% confidence interval, paired $t$-test $p$-value, and the number of folds (out of 5) in which CA-DPP has the lower mean.}
\label{tab:results}
\centering\footnotesize\setlength{\tabcolsep}{1.2pt}
\begin{tabular*}{\textwidth}{@{}l@{\extracolsep{\fill}}cccccccc@{}}
\toprule
 & \multicolumn{4}{c}{$\bar E=0.25$\,mJ/slot} & \multicolumn{4}{c}{$\bar E=2$\,mJ/slot}\\
\cmidrule(lr){2-5}\cmidrule(l){6-9}
Policy & AoII & $\Delta$ (\%) & $p$ & folds & AoII & $\Delta$ (\%) & $p$ & folds\\
\midrule
CA-DPP (proposed) & $0.645\pm0.165$ & -- & -- & -- & $0.390\pm0.085$ & -- & -- & -- \\
AoII threshold~\cite{zakeri2023aoiieh} & $0.925\pm0.258$ & $-22.3\pm11.6$ & $0.002$ & 4/5 & $0.352\pm0.060$ & $+9.5\pm13.2$ & $0.171$ & 3/5 \\
Pre-identification~\cite{gan2025taoi} & $2.556\pm0.284$ & $-74.8\pm6.2$ & $<\!10^{-3}$ & 5/5 & $0.796\pm0.129$ & $-52.0\pm5.2$ & $<\!10^{-3}$ & 5/5 \\
Age threshold~\cite{pehlivanoglu2026orbit} & $2.538\pm0.381$ & $-74.8\pm4.2$ & $<\!10^{-3}$ & 5/5 & $0.784\pm0.114$ & $-51.8\pm5.1$ & $<\!10^{-3}$ & 5/5 \\
Always offload~\cite{basnayaka2025aomi} & $12.451\pm1.288$ & $-94.5\pm1.5$ & $<\!10^{-3}$ & 5/5 & $8.444\pm0.869$ & $-95.2\pm1.1$ & $<\!10^{-3}$ & 5/5 \\
\midrule
CA-DPP, confidence-blind & $0.739\pm0.165$ & $-13.3\pm5.2$ & $0.001$ & 4/5 & $0.345\pm0.056$ & $+10.3\pm9.9$ & $0.051$ & 2/5 \\
CA-DPP, label only & $0.686\pm0.204$ & $+6.7\pm13.8$ & $0.487$ & 1/5 & $0.428\pm0.101$ & $-7.1\pm3.0$ & $0.001$ & 4/5 \\
\midrule
Genie (knows $X$) & \multicolumn{8}{l}{AoII $0.171\pm0.011$, $\Prob$(GS wrong) $0.070$, energy $0.96$\,mJ/slot when unconstrained} \\
\bottomrule
\end{tabular*}
\end{table*}

\subsection{Baselines and Ablations}
No prior policy addresses our exact problem, so we adapt the four closest designs. All baselines observe the same labels, confidences and link states as CA-DPP. They stay within the budget through a token bucket with an initial credit of $0.5$\,J. When a raw upload does not fit in the bucket, they send the label instead if they can.
\begin{itemize}
\item \emph{AoII threshold}, after~\cite{zakeri2023aoiieh}. The source model there has no classifier, so the UAV treats its label as the observed state. It sends the label once the label has disagreed with the GS for $k\in\{1,2,3,5\}$ consecutive slots. This baseline is the strongest one, because it already ignores agreeing labels.
\item \emph{Pre-identification}, after~\cite{gan2025taoi}. The UAV reports every label and uploads the frames labeled flooded, which the ground model re-identifies.
\item \emph{Age threshold}, after~\cite{pehlivanoglu2026orbit}. Once $A\in\{1,3,10\}$ slots have passed since the last delivery, the UAV uploads the frame if $q_2(g)\ge q^\star$, with $q^\star\in\{0.5,0.9\}$ or never, and sends the label otherwise.
\item \emph{Always offload}, after~\cite{basnayaka2025aomi}. The UAV uploads the frame once $N\in\{1,2,5,10,50\}$ slots have passed since the last delivery and never sends labels.
\end{itemize}
Two ablations isolate parts of CA-DPP. The \emph{confidence-blind} version replaces the confidence $c(t)$ by the precision of the label value on tuning data, i.e., by $\Prob\{X=\lab\mid\lab(t)=\lab\}$. Its belief still uses the class prior, and it still distinguishes ``flooded'' labels from ``dry'' ones, but it treats all frames with the same label alike. The \emph{label-only} version never uploads raw data. Finally, a \emph{genie} that knows $X(t)$ and sends an error-free update whenever the GS is wrong gives a floor for the AoII.


\subsection{Evaluation Protocol}\label{res:protocol}
Chips from the same flood event look alike, so results from a single tuning/test split depend strongly on which chips fall in each set. With the official validation/test split, the gains of CA-DPP were much larger than under the cross-validation below, because the two splits differ systematically. We therefore use $5$-fold cross-validation over the $174$ validation and test chips. For each fold, every parameter is chosen on the other four folds. For CA-DPP and its ablations these are $V$ (four values per budget, between $10^{-3}$ and $10$) and the belief update ($\nu\in\{0.25,0.5\}$ with the class prior, or the symmetric update). For the baselines they are the thresholds listed above. The same four folds give the source statistics, the ground error $\ee_{\rm g}(c)$ in seven confidence bins, and the label precisions of the confidence-blind version. Among the parameter values whose energy stays within $5\%$ of the budget, we keep the one with the lowest AoII.

Each selected policy then runs four times on the held-out fold. Each run visits the chips of the fold four times in random order, with fresh fading and blockage, for about $27{,}000$ slots. This gives $20$ held-out runs per policy and setting. Runs with the same index share all random numbers across policies, so the comparison between two policies is paired. We report means with $95\%$ confidence intervals, paired $95\%$ confidence intervals of the relative change in AoII, paired $t$-tests, and the number of folds in which CA-DPP has the lower mean.

\subsection{AoII Under an Energy Budget}\label{res:budget}

\begin{figure*}[t]
\centering
\includegraphics{figures/fig_budget.pdf}
\caption{Performance versus the energy budget $\bar E$: $20$ held-out runs per point ($5$-fold cross-validation over chips) with $95\%$ confidence intervals. (a) Average AoII. (b) Share of slots in which the GS is wrong. (c) AoII reduction of CA-DPP relative to each policy, with paired intervals (positive favors CA-DPP). Baselines adapted from~\cite{zakeri2023aoiieh,gan2025taoi,pehlivanoglu2026orbit,basnayaka2025aomi}.}
\label{fig:budget}
\end{figure*}

Fig.~\ref{fig:budget} and Table~\ref{tab:results} compare CA-DPP with the baselines (Q2). Against the pre-identification and age-threshold policies, CA-DPP lowers the AoII by $75\%$ at $0.25$\,mJ/slot, $52\%$ at $2$\,mJ/slot and $12\%$ at $10$\,mJ/slot, in every fold up to $5$\,mJ/slot ($p<10^{-5}$). The pre-identification policy spends most of its energy on agreeing labels, which are worthless (Proposition~\ref{prop:confirm}). The age-threshold policy reacts to the time since the last delivery, which says little about whether the GS is wrong: an old status of a dry area that is still dry needs no update. Always offloading is an order of magnitude worse, since a raw frame costs $3$ to $21$ labels and often fails far from the GS.

Against the AoII-threshold policy, the outcome depends on the budget. At $0.25$\,mJ/slot, CA-DPP lowers the AoII by $22\%\pm12\%$ ($p=0.002$, four folds of five), and the $99$th percentile of the AoII drops from $19.2$ to $12.3$ slots. Here the UAV can send few labels, and the confidence decides which disagreements deserve them. At $0.5$\,mJ/slot the reduction is $11\%\pm7\%$ and not significant ($p=0.10$). From $1$\,mJ/slot on, the threshold policy with $k=1$ reports every disagreement for about $0.84$\,mJ/slot and matches CA-DPP; the differences range from $-3\%$ to $+14\%$ with $p\ge0.09$. When labels are this cheap, a wrong label is corrected within a slot or two, which the one-step gain of Lemma~\ref{lem:onestep} ignores; a discounted look-ahead did not help.

\begin{figure}[!tb]
\centering
\includegraphics[width=0.8\columnwidth]{figures/fig_path.pdf}
\caption{Sample path of CA-DPP over $120$ slots at $\bar E=0.5$\,mJ/slot ($V=0.01$, $\nu=0.25$). (a) True state (shaded when flooded), GS status, and the slots in which the UAV sends a label (triangles) or uploads the frame (diamond). (b) AoII of the GS status (left axis) and the belief $\mu_t$ that the GS is wrong (right axis). The belief rises with each disagreeing frame until a label resets the AoII.}
\label{fig:path}
\end{figure}

Fig.~\ref{fig:path} shows how CA-DPP works. The AoII spike at the end shows the cost of waiting: the tempered update discounts the first disagreeing frames, so the belief rises more slowly than the true error.

\subsection{Where the Gains Come From}\label{res:gains}
The two ablations answer Q2 and Q3 directly: the confidence-blind version removes the per-frame confidence, and the label-only version removes the raw uploads. At $0.25$\,mJ/slot the per-frame confidence matters: the confidence-blind version has a $15\%$ higher AoII ($p=0.001$). From $1$\,mJ/slot on, however, the confidence-blind version is $6\%$ to $15\%$ better than CA-DPP, and the difference is significant only at $1$\,mJ/slot ($p=0.02$). With this classifier, once energy is not scarce, the precision of each label value gives the belief most of what it needs. A plausible reason is that confident wrong labels come in runs, which the per-frame update over-counts.

Raw uploads pay off at larger budgets. Without them, the AoII stays at $0.43$ from $1$\,mJ/slot on, since the label-only version runs out of useful labels to send and spends only $0.65$\,mJ/slot. CA-DPP spends the rest of the budget on uploads and reaches $0.39$ at $2$\,mJ/slot and $0.34$ at $10$\,mJ/slot ($p\le0.001$). The genie reaches $0.171$ with $0.96$\,mJ/slot; the gap comes mostly from not seeing $X(t)$.

\subsection{Weak Classifiers, Blockage and Payload}\label{res:sens}
\begin{figure*}[t]
\centering
\includegraphics{figures/fig_sensitivity.pdf}
\caption{Sensitivity of the average AoII, $20$ held-out runs per point. (a) On-board accuracy (base model and four noisy variants) and (b) share of blocked slots $\rho_{\rm b}/(\rho_{\rm b}+\rho_{\rm u})$, both at $\bar E=0.5$\,mJ/slot. (c) Raw frame size $L_2$ at $\bar E=2$\,mJ/slot.}
\label{fig:sens}
\end{figure*}

In Fig.~\ref{fig:sens}(a), the AoII of CA-DPP grows only from $0.53$ to $0.64$ as the accuracy drops from $91.8\%$ to $81.8\%$. Over the same range, the AoII of the threshold policy grows from $0.56$ to $1.31$, and that of the confidence-blind version from $0.50$ to $0.94$. With the two noisiest models, it is $27\%$ and $49\%$ better than the threshold policy ($p<10^{-5}$) and $14\%$ and $28\%$ better than the confidence-blind version ($p\le0.003$), in all five folds. The noisiest model rarely says ``flooded'', so over a flooded area it mostly agrees with the wrong ``dry'' status of the GS. The threshold policy, which reacts only to disagreements, spends only $0.32$ of $0.5$\,mJ/slot. CA-DPP spends the full budget: a low-confidence ``dry'' frame barely lowers its belief that the GS is wrong, while a rare confident ``flooded'' frame raises it strongly, since flooded footprints are rare. The confidence-blind belief knows that ``dry'' labels are less reliable on average, but not which ones. Per-frame confidence pays off when the classifier is weak and its errors are lopsided, as is likely for a small model on a low-cost device.

As the share of blocked slots grows from $0$ to $50\%$ (Fig.~\ref{fig:sens}(b)), all policies lose ground and their order is unchanged. CA-DPP stays $9\%$ to $15\%$ below the threshold policy ($0.009\le p\le0.1$). The confidence-blind and label-only versions are $1.5\%$ to $11\%$ below CA-DPP at this budget, so at $0.5$\,mJ/slot raw uploads cost slightly more than they return. Fig.~\ref{fig:sens}(c) varies the size of the raw frame at $2$\,mJ/slot. With $12.5$\,kB frames, CA-DPP is $15\%$ better than the label-only version and matches the threshold policy. With $200$\,kB frames this advantage shrinks to $3.5\%$: uploads cost more and fail more often, and the policy uploads less without retuning.

\subsection{The Weight $V$ and the Role of Mobility}\label{res:V}
\begin{figure}[!tb]
\centering
\includegraphics[width=0.8\columnwidth]{figures/fig_V_mix.pdf}
\caption{(a) Average AoII of CA-DPP (left axis) and energy spent (right axis, dashed) versus $V$ at $\bar E=0.5$\,mJ/slot. In the shaded range the energy stays within $4\%$ of the budget (dotted); a larger $V$ lowers the AoII but overspends. (b) Share of unblocked slots with a label or a raw upload versus the distance to the GS, at $\bar E=2$\,mJ/slot. Uploads dominate near the GS; labels take over beyond about $0.8$\,km.}
\label{fig:Vmix}
\end{figure}

Fig.~\ref{fig:Vmix}(a) shows the trade-off governed by $V$ at $0.5$\,mJ/slot. Raising $V$ from $10^{-3}$ to $3$ lowers the AoII from $0.79$ to $0.41$. A larger $V$, however, lets the debt grow before it raises the price: over about $27{,}000$ slots the energy spent rises from $0.50$ to $1.51$\,mJ/slot. Values of $V$ up to $10^{-2}$ stay within $4\%$ of the budget. Theorem~\ref{thm:budget} holds for every $V$, but the transient excess grows with $V$, so $V$ must match the mission length.

Fig.~\ref{fig:Vmix}(b) shows how CA-DPP exploits the trajectory. Within $850$\,m of the GS, where uploads are cheap and reliable, it uploads the frame in $5.9\%$ of the unblocked slots and sends labels in $2.0\%$. Beyond $2.15$\,km it uploads in only $0.1\%$ and sends labels in $3.8\%$. The UAV asks the ground when asking is cheap, although the distance enters~\eqref{eq:rule} only through $q_a(g)$ and $E_a(g)$.
\subsection{Distance to the Optimal Policy}\label{res:opt}
Remark~\ref{rem:opt} leaves the gap between CA-DPP and the optimal policy of (P) open, and computing that optimum is out of reach, since the belief has $\Delta_{\max}+1$ entries. A lower bound on it can be computed, however. A controller that also observes $X(t)$, and therefore the AoII and whether the on-board label is right, can imitate any policy of (P), so the optimal AoII of this full-information controller under the same budget is a lower bound on the optimum of (P). Its problem is a constrained Markov decision process over $X(t)$, the AoII, the orbit position and the blockage state. We solve its Lagrangian relaxation by value iteration for a range of energy prices, with the AoII capped (which can only lower the bound), and obtain the bound at each energy level by weak duality. To compare against the bound without confounding factors, we use a synthetic model that satisfies Assumption~\ref{as:inf}: $X(t)$ follows the Markov chain fitted to the data, and in each slot the outputs of the two classifiers are drawn jointly from data slots with the same state, so the errors are independent over time. CA-DPP then uses $\nu=1$. CA-DPP and the AoII-threshold policy are tuned as in Section~\ref{res:protocol} and run six times for $10^5$ slots.

\begin{table}[t]
\caption{Synthetic model with independent errors: average AoII of CA-DPP and of the AoII-threshold policy (mean of six runs), and the lower bound at the energy each policy spends. The gap is the excess over the bound; the true gap to the optimal policy is smaller.}\label{tab:optgap}
\centering\footnotesize
\begin{tabular}{@{}ccccccc@{}}
\toprule
 & \multicolumn{3}{c}{CA-DPP} & \multicolumn{3}{c}{AoII threshold}\\
\cmidrule(lr){2-4}\cmidrule(l){5-7}
$\bar E$ (mJ/slot) & AoII & bound & gap & AoII & bound & gap\\
\midrule
$0.25$ & $0.878$ & $0.701$ & $+25\%$ & $1.208$ & $0.706$ & $+71\%$\\
$0.5$ & $0.624$ & $0.497$ & $+25\%$ & $0.941$ & $0.686$ & $+37\%$\\
$1$ & $0.389$ & $0.320$ & $+22\%$ & $0.464$ & $0.335$ & $+38\%$\\
$2$ & $0.348$ & $0.306$ & $+14\%$ & $0.464$ & $0.335$ & $+38\%$\\
\bottomrule
\end{tabular}
\end{table}

Table~\ref{tab:optgap} shows that CA-DPP is within $14\%$ to $25\%$ of the bound at every budget, and the gap shrinks as energy grows. The bound credits the controller with knowing $X(t)$, which no policy of (P) can match, so CA-DPP is at least this close to the optimal policy. The AoII-threshold policy is $37\%$ to $71\%$ above the bound; at the larger budgets it spends less than allowed, because even with $k=1$ it sends only labels that disagree with the GS. In this model CA-DPP also stays ahead of the threshold policy at $2$\,mJ/slot, by $25\%$ in AoII, while on the real traces the two tie at this budget (Section~\ref{res:budget}). The difference between the two settings is the correlation of the errors, which the next subsection examines.

\subsection{Correlated Errors}\label{res:corr}
On the real traces the on-board errors come in runs (Section~\ref{res:cls}). To isolate their effect, we extend the synthetic model with an error indicator that follows a two-state Markov chain with the measured error rate $8.2\%$ and persistence $\Prob\{\text{wrong at }t+1\mid\text{wrong at }t\}=\varrho$. Given the state and the indicator, the confidence and the ground output are drawn from data slots with the same state and the same correctness, so the classifier stays calibrated. We vary $\varrho$ from $0.082$ (independent errors) to $0.8$; the real traces have $\varrho=0.56$. CA-DPP runs with the rule of Section~\ref{sec:update} and, for comparison, with $\nu=1$.

\begin{table}[t]
\caption{Synthetic model with correlated errors: AoII reduction of CA-DPP relative to the AoII-threshold policy (mean of six paired runs; negative favors CA-DPP), with $\nu$ from the rule and with $\nu=1$.}\label{tab:corr}
\centering\footnotesize
\begin{tabular}{@{}cccccc@{}}
\toprule
 & & \multicolumn{2}{c}{$\bar E=0.5$\,mJ/slot} & \multicolumn{2}{c}{$\bar E=2$\,mJ/slot}\\
\cmidrule(lr){3-4}\cmidrule(l){5-6}
$\varrho$ & rule $\nu$ & rule $\nu$ & $\nu=1$ & rule $\nu$ & $\nu=1$\\
\midrule
$0.082$ & $1.00$ & $-23\%$ & $-23\%$ & $-23\%$ & $-23\%$\\
$0.3$ & $0.62$ & $-21\%$ & $-21\%$ & $-7\%$ & $-6\%$\\
$0.56$ & $0.32$ & $-20\%$ & $-17\%$ & $-4\%$ & $-1\%$\\
$0.8$ & $0.12$ & $-34\%$ & $-13\%$ & $-10\%$ & $-1\%$\\
\bottomrule
\end{tabular}
\end{table}

Table~\ref{tab:corr} shows two effects. With independent errors, CA-DPP lowers the AoII of the threshold policy by about $23\%$ at both budgets. As the errors become more persistent, the advantage at $2$\,mJ/slot shrinks, because a run of wrong labels looks like strong evidence to a belief that counts each frame separately; with $\nu=1$ it disappears. The rule for $\nu$ recovers part of it, and at $0.5$\,mJ/slot it keeps CA-DPP ahead at every persistence. This suggests that the tie on the real traces at large budgets comes from the correlated errors, which the belief does not model, and a belief that models the error process (Section~\ref{sec:discussion}) is the next step.

\subsection{Robustness Checks}\label{res:robust}
\emph{Choice of $\nu$.} The main results select $\nu$ on the tuning folds. Table~\ref{tab:nu} repeats the trace-driven evaluation with $\nu$ set instead by the correlation rule of Section~\ref{sec:update}, which gives $\nu$ between $0.31$ and $0.35$ for the base model and between $0.30$ and $0.42$ for the noisy variants; everything else, including the folds and the random numbers, is unchanged. With the rule, CA-DPP lowers the AoII of the threshold policy by $16\%$ (paired relative change) at $0.5$\,mJ/slot (paired $p=0.003$) and by $11\%$ and $15\%$ at $5$ and $10$\,mJ/slot ($p=0.015$ and $p=0.008$ on the five fold means), where the tuned version ties. At $0.25$\,mJ/slot, however, its AoII rises from $0.65$ to $0.87$ and it only ties with the threshold policy, and with the weakest classifier its AoII rises from $0.64$ to $0.78$, still $41\%$ below the threshold policy. The best discount therefore depends on the budget, and choosing it per budget on tuning data, as in the main results, is the safer default. Repeating the rule run with ten folds gives the same picture: a tie at $0.25$\,mJ/slot and reductions of $7\%$ at $0.5$\,mJ/slot and of $31\%$ and $40\%$ with the two weakest classifiers.

\begin{table}[t]
\caption{Average AoII of CA-DPP with $\nu$ selected on the tuning folds (main results) and with $\nu$ from the correlation rule, and of the AoII-threshold policy (mean of $20$ held-out runs, $5$-fold cross-validation). The last two rows use the two weakest classifiers at $\bar E=0.5$\,mJ/slot.}\label{tab:nu}
\centering\footnotesize
\begin{tabular}{@{}lccc@{}}
\toprule
 & \multicolumn{2}{c}{CA-DPP} & AoII\\
\cmidrule(lr){2-3}
Setting & tuned $\nu$ & rule $\nu$ & threshold\\
\midrule
$\bar E=0.25$\,mJ/slot & $0.645$ & $0.867$ & $0.925$\\
$\bar E=0.5$\,mJ/slot & $0.526$ & $0.479$ & $0.564$\\
$\bar E=1$\,mJ/slot & $0.418$ & $0.356$ & $0.363$\\
$\bar E=2$\,mJ/slot & $0.390$ & $0.337$ & $0.352$\\
$\bar E=5$\,mJ/slot & $0.360$ & $0.309$ & $0.352$\\
$\bar E=10$\,mJ/slot & $0.341$ & $0.291$ & $0.352$\\
accuracy $83.9\%$ & $0.581$ & $0.570$ & $0.809$\\
accuracy $81.8\%$ & $0.643$ & $0.775$ & $1.312$\\
\bottomrule
\end{tabular}
\end{table}

The selection itself is noisy. With $\nu$ fixed and only $V$ tuned, CA-DPP reaches an AoII of $0.41$, $0.44$, $0.47$, $0.55$ and $0.72$ at $0.5$\,mJ/slot for $\nu=0.14$, $0.25$, $0.32$, $0.5$ and $1$, and $0.34$, $0.34$, $0.34$, $0.36$ and $0.41$ at $2$\,mJ/slot. Exact Bayes' rule ($\nu=1$) is clearly the worst, which confirms that the tempering matters, and several fixed values do better than the tuned version in Table~\ref{tab:results}, because the selection on four folds sometimes picks a poor value. The main results are therefore conservative.

\emph{Exact energy.} The energy model~\eqref{eq:energy} replaces the fading average by its value at the mean SNR. We repeated the budget and accuracy sweeps with the exact average over the fading, computed by numerical integration, and retuned every policy. The AoII of CA-DPP changes by at most $4.2\%$ across budgets and $6.7\%$ across classifiers, its relative change against the threshold policy moves by at most $6$ percentage points, and no significant comparison changes sign.

\emph{A belief-threshold baseline.} Once the belief is available, a simpler policy may seem enough. We therefore tested a rule that sends the label whenever the belief $\mu_t$ that the GS is wrong exceeds a threshold. It uses the same belief as CA-DPP with $\nu=0.25$, never uploads raw frames, and has its threshold chosen per budget on the tuning folds. At $0.25$, $0.5$, $1$, $2$, $5$ and $10$\,mJ/slot its average AoII is $0.75$, $0.47$, $0.37$, $0.37$, $0.37$ and $0.37$, against $0.65$, $0.53$, $0.42$, $0.39$, $0.36$ and $0.34$ for CA-DPP, and none of these differences is significant (paired $p\ge0.10$). The belief is therefore the part of CA-DPP that carries most of its gain over the AoII-threshold policy. What the controller adds is control of the energy. The threshold rule meets a budget only through tuning on past data, with no guarantee, and its energy use stops near $0.65$\,mJ/slot because it never uploads; CA-DPP tracks the budget through the queue with the guarantee of Theorem~\ref{thm:budget} and adapts its thresholds to the link and the energy price in every slot (Proposition~\ref{prop:threshold}).

\emph{Statistics.} Runs on the same fold share chips, so the $20$ held-out runs are not fully independent. Tested on the five fold means instead, the reductions with the two weakest classifiers remain significant ($p=0.035$ and $p=0.014$, all five folds), whereas the $0.25$\,mJ/slot reduction, seen in four of five folds, has $p=0.17$. The large gains over the pre-identification and age-threshold policies hold in every fold up to $5$\,mJ/slot and in four of five folds at $10$\,mJ/slot.

%% =====================================================================
\section{Discussion}\label{sec:discussion}
%% =====================================================================
\subsection{Answers to the Research Questions}
\emph{Q1.} Never send an agreeing label; otherwise take the action whose value $q_a(g)G_a$ most exceeds its priced energy, which gives confidence thresholds ordered silent, raw upload, label (Proposition~\ref{prop:threshold}), and the price keeps the budget (Theorem~\ref{thm:budget}).

\emph{Q2.} The per-frame confidence helps when updates are scarce, either because energy is scarce or because the classifier is weak. The UAV must then choose which few disagreements to report, and the confidence ranks them. When energy is ample and the classifier accurate, every disagreement can be reported and a simple rule does as well.

\emph{Q3.} The value of raw uploads depends on the link more than on the classifier. Uploads pay off at budgets of $2$\,mJ/slot and above, near the GS and with small frames.

Two design lessons follow. First, never send a label that agrees with the receiver (Proposition~\ref{prop:confirm}); much of the $52\%$ to $75\%$ gap to the pre-identification and age-threshold policies comes from the energy they spend on such labels. Second, calibrate the classifier on board: isotonic calibration is a lookup table, and the prior correction~\eqref{eq:lik} matters because monitored events are rare.

\subsection{Correlated Errors}
The main weakness of CA-DPP at large budgets is Assumption~\ref{as:inf}(ii): the on-board errors come in runs (Section~\ref{res:cls}). The tempered update~\eqref{eq:post} discounts every frame equally and cannot tell new evidence from a repeated error. Section~\ref{res:corr} shows that this costs CA-DPP most of its advantage over the threshold policy at large budgets, and that the rule for $\nu$ recovers part of it. A hidden Markov model of the classifier state could recover more, at twice the belief size.

\subsection{Other IoT Settings}
Nothing in the model is specific to UAVs or floods. The link and energy models fit any device whose channel varies predictably, such as a low-Earth-orbit satellite that sees a ground station only during passes. A non-binary scene needs a belief over the AoII and the current class, but the gains keep their structure: a right update removes AoII and a wrong one creates it.

\subsection{Limitations}
Five limitations remain. First, the policy is myopic; its gap to the optimal policy of (P) is bounded only numerically, in a model with independent errors (Section~\ref{res:opt}). Second, Proposition~\ref{prop:threshold} holds for the per-slot rule~\eqref{eq:rule} within each confidence bin; whether the optimal policy of (P) shares it is open. Third, Sentinel-2 imagery stands in for UAV frames (FLAME~\cite{shamsoshoara2021flame} and AIDER~\cite{kyrkou2020aider} are UAV alternatives). Fourth, the energy model is analytic; replacing the mean-SNR approximation~\eqref{eq:energy} by the exact fading average changes none of the conclusions (Section~\ref{res:robust}), but measured energy is still missing. Fifth, the imagery is not tied to the UAV position, so route planning is left out.

%% =====================================================================
\section{Conclusion}\label{sec:concl}
%% =====================================================================
We studied status updates from an IoT device whose sensor is its own classifier, with a UAV mapping a flood as the running example. A belief over the AoII of the GS status, updated with the calibrated confidence of each frame, gives closed-form values of the three actions: agreeing labels are worthless, and the per-slot decision has confidence thresholds ordered silent, raw, label. CA-DPP, built on them, meets the energy budget on every sample path. On flood-mapping traces under chip-level cross-validation, it beats adaptations of three recent designs at every budget up to $5$\,mJ per slot, and an AoII-threshold policy when energy is scarce or the classifier is weak. Raw uploads pay off near the GS at larger budgets. Next steps are a belief that models correlated errors, and UAV imagery with measured energy.

\section*{Declaration of competing interest}
The author declares no known competing financial interests or personal relationships that could have appeared to influence the work reported in this paper.

\section*{Funding}
This research did not receive any specific grant from funding agencies in the public, commercial, or not-for-profit sectors.

\section*{Data availability}
The Sen1Floods11 dataset~\cite{bonafilia2020sen1floods11} is publicly available. The simulation code, inference traces and scripts that produce every figure and table are available at \url{https://github.com/sshawky78AUA/confidence-aware-status-update-IoT}.

\section*{Declaration of generative AI and AI-assisted technologies in the writing process}
During the preparation of this work the author used Claude (Anthropic) in order to assist with drafting and editing the text, writing the simulation and plotting code, and preparing the figures. After using this tool, the author reviewed and edited the content as needed and takes full responsibility for the content of the publication.

\appendix
\section{Proof of Lemma~\ref{lem:mono}}\label{app:mono}
Let $y=\Lambda_t^{\nu}$, which increases with $\Lambda_t$. By~\eqref{eq:post}, the mass at every $\delta>0$ is multiplied by $y/D$ and the mass at $\delta=0$ by $1/D$, with $D=1-\mu^-_t+\mu^-_ty$. Hence
\[
S_t=\frac{y\,S^-_t}{D},\qquad 1-\mu_t=\frac{1-\mu^-_t}{D},
\]
and $AS_t-B(1-\mu_t)=N/D$ with $N=AyS^-_t-B(1-\mu^-_t)$. Both $N$ and $D$ are affine in $y$, so the derivative of $N/D$ with respect to $y$ has the sign of
\begin{multline*}
AS^-_t\big(1-\mu^-_t+\mu^-_ty\big)-\big(AyS^-_t-B(1-\mu^-_t)\big)\mu^-_t\\
=(1-\mu^-_t)\big(AS^-_t+B\mu^-_t\big)\ \ge\ 0.
\end{multline*}
Since $D>0$, $N/D$ is nondecreasing in $y$, and hence in $\Lambda_t$. \hfill$\blacksquare$

\section{Proof of Proposition~\ref{prop:threshold}}\label{app:threshold}
Let $f_a=Vq_aG_a-\lambda E_a$ be the objective of~\eqref{eq:rule}, so $f_0=0$. Since $r,w<\frac12$ and $\ee_{\rm g}\le1$, both gains have the form $AS_t-B(1-\mu_t)$ with $A,B\ge0$: for $G_2$, $A=(1-\ee_{\rm g})(1-2w)$ and $B=\ee_{\rm g}(1-2r)$, and for $G_1$ with $d(t)=1$, $A=1-2w$ and $B=1-2r$. The difference $q_1G_1-q_2G_2$ has the same form, with
\[
A=(1-2w)\big(q_1-q_2(1-\ee_{\rm g})\big),\qquad B=(1-2r)\big(q_1-q_2\ee_{\rm g}\big),
\]
both nonnegative because $q_2\le q_1$. By Lemma~\ref{lem:mono}, $f_2-f_0=Vq_2G_2-\lambda E_2$ and $f_1-f_2=V(q_1G_1-q_2G_2)-\lambda(E_1-E_2)$ are nondecreasing in $\Lambda_t$, and so is their sum $f_1-f_0$.

(a) If $d(t)=1$, then $\Lambda_t$ increases in $c$. Order the actions as $0\prec2\prec1$. The differences between consecutive actions in this order are nondecreasing in $c$, so $f_a$ has increasing differences in $(a,c)$. By Topkis' theorem, the set of maximizers of $f_a$ is nondecreasing in $c$ in the strong set order, so its smallest and largest elements are both nondecreasing in $c$. Algorithm~\ref{alg:dpp} selects the smallest maximizer, breaking ties toward silence and then toward the raw upload. Hence the action moves from 0 to 2 to 1 as $c$ grows, which defines $c_{\rm r}\le c_{\lab}$.

(b) If $d(t)=0$, then $G_1=0$ by Proposition~\ref{prop:confirm}, so $f_1=-\lambda E_1\le f_0$ and the label is never strictly preferred to silence. Now $\Lambda_t$ decreases in $c$, so $f_2-f_0$ is nonincreasing in $c$, and the set of $c$ for which $f_2>f_0$ is an interval of the form $\{c\le c^0_{\rm r}\}$.

(c) For fixed $\mu^-_t$ and $\Lambda_t$, $S^-_t$ appears in $N/D$ only in the numerator, with coefficient $Ay\ge0$. So $f_a-f_0$ is nondecreasing in $S^-_t$ for $a\in\{1,2\}$, and a slot in which some action beats silence still has one when $S^-_t$ grows. Each $f_a-f_0$ with $a\neq0$ decreases in $\lambda$, so a slot in which silence is optimal remains so when $\lambda$ grows. \hfill$\blacksquare$

\section{Proof of Theorem~\ref{thm:budget}}\label{app:budget}
We first bound the gains. Since the belief is truncated at $\Delta_{\max}$, $S_t=\E_{\pi_t}[(\Delta+1)\ind{\Delta>0}]\le\Delta_{\max}+1$. By~\eqref{eq:G1} and~\eqref{eq:G2}, $G_a\le S_t$ for $a\in\{1,2\}$, because the negative terms can be dropped and $(1-2w)(1-\ee_{\rm g})\le1$. With $q_a\le1$, $Vq_aG_a\le V(\Delta_{\max}+1)$.

Let $z^\star=V(\Delta_{\max}+1)/E_{\min}$. If $Z(t)>z^\star$, then for $a\in\{1,2\}$,
\[
Vq_aG_a-Z(t)E_a(g)<V(\Delta_{\max}+1)-z^\star E_{\min}=0=f_0,
\]
so CA-DPP stays silent, and~\eqref{eq:Z} gives $Z(t+1)=\max\{Z(t)-\bar E,0\}\le Z(t)$. If $Z(t)\le z^\star$, then $Z(t+1)\le Z(t)+E_{\max}\le z^\star+E_{\max}$. Starting from $Z(0)=0$, induction on $t$ gives $Z(t)\le z^\star+E_{\max}=Z_{\max}$ for all $t$. Inserting this bound in~\eqref{eq:Zbound} gives the energy bound for every $T$, and taking expectations and $T\to\infty$ gives~\eqref{eq:budget}. \hfill$\blacksquare$

\setlength{\bibsep}{0pt}
\bibliographystyle{elsarticle-num}
\bibliography{refs}

\end{document}
